% =====================================================================
% Methods and Design Companion to "Degrees of Separation" (main.tex).
% A technical retrospective of the STATIC analysis.
% Preamble reused verbatim from main.tex so the two compile in one style.
% NOTE: no LaTeX toolchain is installed in the authoring environment
% (pdflatex/bibtex absent); compile externally with
%   pdflatex methods_companion && bibtex methods_companion && pdflatex x2.
% =====================================================================
\documentclass[11pt]{article}
\usepackage[margin=1in]{geometry}
\usepackage{graphicx}
\usepackage{booktabs}
\usepackage{amsmath,amssymb}
\usepackage{natbib}
\usepackage{hyperref}
\usepackage{xcolor}
\usepackage{enumitem}
\graphicspath{{../outputs/figures/}{../results/figures/}}
\newcommand{\todo}[1]{\textcolor{red}{[TODO: #1]}}
\newcommand{\coauthor}[1]{\textcolor{blue!60!black}{[COAUTHOR: #1]}}

\title{\textbf{Degrees of Separation --- Methods and Design Companion}\\[4pt]
\large A technical retrospective of the static analysis:\\
every design decision, every equation, every number}
\author{Sora Yng \and \coauthor{authorship/order to be set; see COAUTHOR\_NOTES.md}}
\date{Companion to \texttt{main.tex}, assembled from the analysis record. Not for circulation.}

\begin{document}
\maketitle

\begin{abstract}
This is the methods-and-design companion to ``Degrees of Separation: a measurable wedge between academic
prestige and the labour-market pricing of fields'' (\texttt{main.tex}). It is not a second paper and it is not
a future-work agenda; it is a \emph{technical retrospective} of the static analysis. Where the paper
\emph{states} a result, this document \emph{explains} it: every estimator is presented with a ``Reading.''
that defines its symbols, justifies its functional form against the alternatives, and states what it buys and
what it assumes; and every number that appears in the paper's claims ledger (\texttt{CLAIMS\_LEDGER.md}) is not
merely reported but \emph{decoded} --- why it takes the value it does, and what that value licenses and
forecloses. The companion follows the analysis pipeline in order (data construction; the gap definition and its
reliability; the discipline decomposition; the licensing channel; the ruled-out alternatives; the two-signal
model; the robustness battery; the temporal and applied extensions) and then turns, as a first-class matter, to
the single deepest interpretive constraint on the whole construction --- the \emph{measurement asymmetry}
between a richly relational academic axis and a scalar earnings axis (\S\ref{sec:asymmetry}). All statistics are
copied from the same \texttt{*\_RESULT.md} record the paper draws on; nothing here is recomputed or invented.
\end{abstract}

% =====================================================================
\section{What this document is, and the interpretive throughline}\label{sec:companion-intro}

\paragraph{Purpose and stance.} The main paper compresses a long analysis into a small number of declarative
results. That compression is appropriate for a paper but it hides the part a methodologically serious reader
most wants to see: \emph{why} each measurement and estimator was chosen over its alternatives, and \emph{why}
each headline number takes the precise value it does rather than some other. This companion restores that
layer. It is organised so that it can be read alongside the paper section by section, and it holds itself to
two hard standards. First, no equation appears without a ``Reading.'' paragraph that defines every symbol,
explains the choice of that functional form or estimator over the competitors, and is explicit about what the
choice buys and what it assumes. Second, no number is left un-decoded: every value in the paper's claims ledger
is presented \emph{and} interpreted --- a reproduction correlation below one is diagnosed as a data ceiling
rather than a bug; an intraclass correlation and a one-way $\eta^2$ that disagree are reconciled by the
shrinkage that separates them; a null coefficient is certified as real by the near-orthogonality of its
regressors; a residual ranking that is unchanged to three decimals is read as redundancy rather than mechanism.
The document is deliberately exhaustive; completeness and interpretation depth, not brevity, are its point.

\paragraph{The central object, restated.} Everything below concerns one statistic. For a field $f$ observed
across the institutions $i$ that teach it,
\[
\mathrm{gap}_f \;=\; 1 - \mathrm{Spearman}_i\!\left(\mathrm{prestige}_{i,f},\ \mathrm{earnings}_{i,f}\right),
\]
the within-field, across-institution rank disagreement between an academic-reputation (AR) ordering and an
employer-reputation (ER) ordering of the same institution-by-field cells. The paper reads a large gap as a
field in which academic standing and market pay are \emph{decoupled} and a small gap as one in which they are
\emph{integrated}. Its symbols, and the rationale for every design choice in it --- why a rank statistic, why the $1-(\cdot)$ orientation, why within-field across-institution --- are decoded in full in \S\ref{sec:companion-gap}; here it serves only to fix the object. The companion's job is to make every step from that definition to the paper's claims
inspectable.

\paragraph{The throughline.} One structural fact organises the whole retrospective and recurs at every stage:
the two sides of the gap are measured with profoundly \emph{asymmetric richness}. The academic axis is a
relational network object --- a continuous SpringRank score distilled from thousands of faculty-hiring
decisions --- while the employer axis, as operationalised, is a single scalar, the median earnings of one
slice of graduates at one career horizon. This asymmetry is not a blemish to be apologised for in a footnote;
it is the binding interpretive constraint on what the gap and, especially, its irreducible residual can mean.
We flag it as we go and then confront it directly and substantively in \S\ref{sec:asymmetry}, which is the
argumentative centre of gravity of this companion.

\paragraph{A note on cross-references.} A cross-reference written as ``\S$N$'' (for a bare number $N$, e.g.\
\S4 or \S\S4--5) points to the correspondingly numbered section of the \emph{main paper}, \texttt{main.tex}
(framework \S2, data \S3, the gap \S4, channels \S5, the model \S6, robustness \S7, discussion \S8). Sections of
\emph{this} companion are always referred to by title or by an explicit ``\S\ref{sec:asymmetry}''-style label,
never by a bare number, so the two numbering systems do not collide.

\paragraph{Roadmap.} Section~\ref{sec:companion-data} is the construction log for the two axes and decodes the
coverage and reproduction numbers. Section~\ref{sec:companion-gap} sets out the gap estimator, the failure mode
it is exposed to, the diagnostic battery, and the reliability filter. Section~\ref{sec:companion-decomp} decodes
the headline multilevel decomposition and the integrated--decoupled typology.
Section~\ref{sec:companion-channels} decodes the licensing channel and the redundancy of the additional
channels. Section~\ref{sec:companion-ruleout} decodes the three ruled-out alternatives as scope conditions.
Section~\ref{sec:model-companion} presents the two-signal model and is explicit about its unvalidated status.
Section~\ref{sec:robust-companion} decodes the robustness battery. Section~\ref{sec:temporal-applied} decodes the
temporal revaluation and the applied extensions. Section~\ref{sec:asymmetry} develops the measurement asymmetry
as a first-class methodological argument, and Section~\ref{sec:closing} states tightly what the static paper
establishes and what remains open --- all of it on the ER side.
\section{Data construction and its numbers}
\label{sec:companion-data}

This section is the construction log behind \S3 of the main paper. The paper \emph{states} that the gap is built from a continuous SpringRank reconstruction of the faculty-hiring network on one axis and median field-of-study earnings on the other, joined by CIP and cross-validated against PSEO. Here we explain \emph{how} each axis is built, \emph{why} each construction choice was made over its alternatives, and \emph{what} each coverage and reproduction number certifies. The organising principle is that the gap $\mathrm{gap}_f = 1 - \mathrm{Spearman}_i(\text{prestige}_{i,f}, \text{earnings}_{i,f})$ is a within-field rank statistic, so the data only has to support \emph{rankings} of institutions inside a field --- but two downstream uses (the licensing blend of \S5 and the applied CS ranking of \S8) need score \emph{differences}, not just orderings, and that single requirement drives the central construction decision on the academic axis.

\subsection{Two axes, two valuation systems: why AR and ER are measured separately}

\paragraph{The measurement decision.} The object of study is a wedge between two prices for the same human-capital bundle (institution $\times$ field): an \emph{academic reputation} (AR) axis realised in who hires whose PhDs, and an \emph{employer reputation} (ER) axis realised in earnings. The deliberate decision is to measure these on two structurally independent data systems --- the faculty-hiring network of \citet{wapman2022quantifying} for AR, and College Scorecard field-of-study earnings \citep{scorecard} for ER --- rather than from a single ratings instrument. The payoff is that the gap cannot be an artifact of a shared measurement frame: a survey that asked respondents to rate ``prestige'' and ``career outcomes'' would correlate the two axes through common-method variance, manufacturing agreement (a small gap) wherever respondents conflate the constructs. Two unrelated administrative systems --- a hiring census and a tax/earnings record --- have no such shared error channel, so an observed gap is a genuine disagreement between the academy and the labour market, not an instrument echo. The cost is that the two systems do not align perfectly in coverage; reconciling them is the entire content of the coverage numbers below. The systematic \emph{asymmetry} between the AR and ER axes --- the fact that they fail in different directions and on different institutions --- is a result in its own right and is treated in its dedicated section; here we only establish how each axis is built and what its coverage certifies.

\subsection{The academic axis: why recompute a continuous SpringRank}

\paragraph{The released ranks are ordinal; the analysis needs cardinal.} The Wapman deposit ships per-field \emph{ordinal} prestige ranks (and a $[0,1]$-scaled \texttt{PrestigeRank}), but no analysis code and no documented hyperparameters (the deposit inventory in \texttt{GENREG\_RECALIBRATION\_RESULT.md} lists the SpringRank code as ``NOT FOUND'' and $\alpha$/preprocessing as ``NOT DOCUMENTED''). For the gap itself, ordinal ranks suffice --- Spearman is rank-invariant. But the licensing blend in \S5 and the applied CS ranking in \S8 form weighted combinations of prestige and earnings, and a weighted average of two \emph{orderings} is meaningless: the spacing between rank 1 and rank 2 must be commensurate with the spacing between rank 50 and rank 51 for a blend to be defined. That requires a \emph{cardinal} score whose differences carry information. We therefore recompute the continuous SpringRank score \citep{debacco2018physical} directly from the public field-level hiring edges.

\paragraph{What SpringRank is.} SpringRank treats the directed hiring network as a physical system of springs. Let $A$ be the weighted adjacency matrix of the hiring network for a field: $A_{ij}$ is the number of faculty placements from PhD-granting institution $i$ into hiring institution $j$ (an edge ``points'' from where a scholar trained to where they were hired). Each node $i$ is assigned a real-valued rank $s_i$, and every directed edge is modelled as a spring of rest length $1$ connecting $s_i$ to $s_j$. The scores are the configuration that minimises the total elastic energy:
\[
H(s) \;=\; \tfrac{1}{2}\sum_{i,j} A_{ij}\,\bigl(s_i - s_j - 1\bigr)^2 \;+\; \tfrac{\alpha}{2}\sum_i s_i^2,
\qquad
s^\star = \arg\min_s H(s).
\]
Because $H$ is quadratic in $s$, the minimiser solves a single sparse linear system. Writing $d^{\text{out}}_i = \sum_j A_{ij}$, $d^{\text{in}}_i = \sum_j A_{ji}$, $D = \operatorname{diag}(d^{\text{out}} + d^{\text{in}})$, and using the regularisation $\alpha$ (a spring tethering every node to the origin so the system is non-degenerate), the first-order condition $\nabla H = 0$ is the linear system
\[
\bigl(\,D - (A + A^{\!\top}) + \alpha I\,\bigr)\, s^\star \;=\; \bigl(d^{\text{out}} - d^{\text{in}}\bigr).
\]

\paragraph{Reading.} Here $s_i$ is the continuous prestige score of institution $i$ (higher = more prestigious; it is the cardinal object our blend needs), $A_{ij}$ the count of $i\to j$ faculty placements, $\alpha>0$ a regularisation/tethering strength, $I$ the identity, and the right-hand side $d^{\text{out}}_i - d^{\text{in}}_i$ is institution $i$'s \emph{net} placement flow --- how many more of its PhDs it sends out as faculty than it takes in. The minimised energy $H$ rewards a configuration in which every hiring edge spans exactly one ``rung'' ($s_i - s_j \approx 1$): a placement from a high node to a low node is consistent (low energy), an upward placement is anomalous (high energy), so the optimum is the ranking that makes the observed hiring flow look maximally like a top-down cascade. \emph{Why this estimator over the alternatives:} a pure ordinal sort (or PageRank/eigenvector centrality) returns only an ordering and would force us back to meaningless rank-arithmetic in the blend; minimum-violation rank aggregation is combinatorial and unstable on sparse field subgraphs; SpringRank instead returns a real-valued, regularised score as the unique solution of a \emph{linear} system, so it is cheap, deterministic, and --- critically --- its score \emph{differences} are physically interpretable as the number of hiring rungs between two institutions. \emph{What it buys:} cardinal spacing (so $s_i - s_j$ is meaningful in the licensing blend and the CS ranking) plus exact reproducibility from public edges. \emph{What it assumes:} that net placement flow encodes a single latent prestige dimension and that the rest-length-$1$ spring is the right local model of a hiring step --- i.e. that one ``hire'' is one unit of prestige distance, with $\alpha$ controlling how hard nodes with little data are pulled toward the centre. The regulariser $\alpha$ is a genuine free parameter; the recalibration sweep (next subsection) selects it on reproduction rather than leaving it implicit.

\subsection{What the from-scratch reproduction number certifies}

\paragraph{$\rho \approx 0.77$ is a data ceiling, not a code bug.} Recomputing SpringRank from the public field edges does not reproduce Wapman's published ranks exactly. Sweeping the canonical \texttt{cdebacco/LarremoreLab} implementation over $\alpha \in \{0, 0.1, 0.5, 1, 2\}$ and binary-vs-raw edge weights, the field-averaged Spearman against the published ranks peaks at $\rho = 0.772$ at $\alpha = 0.5$ with binary weights, and the whole surface tops out near $0.80$ (per-field range $0.72$--$0.84$):
\[
\max_{\alpha,\,\text{weighting}} \;\overline{\mathrm{Spearman}}\bigl(s^\star(\text{public edges}),\ \text{rank}^{\text{Wapman}}\bigr) \;=\; 0.772 .
\]

\paragraph{Reading.} The left-hand side is the best reproduction the public data permits, maximised over the only two knobs the estimator exposes ($\alpha$ and binarisation). The value is decisive precisely because it is a \emph{maximum} and it still falls short of $1.0$: neither regularisation strength nor edge-weighting can push past $\approx 0.80$. This rules out the two ``code-bug'' explanations --- a mis-set hyperparameter (the sweep covers it) or a weighting error (binary vs raw both tested) --- and leaves only the data explanation. Wapman ranked the \emph{proprietary} AARC faculty census; the public deposit ships \emph{aggregated} edge lists that are lossy relative to that census (counts are coarsened, low-frequency edges are thinned). A SpringRank fit to a lossy projection of the true network cannot recover the ranks computed on the full network, no matter how well-tuned. Hence $\rho \approx 0.77$ is a \emph{ceiling imposed by the input data}, not a defect of our pipeline; $0.80$ is simply how much of the proprietary ordering survives in the public edges. This number does double duty: it is the honest fidelity statement for our AR axis, and --- because re-estimating ranks on the public edges each bootstrap draw would inject this irreducible $\approx 0.2$ of data-mismatch noise into every replicate --- it is the justification for the \emph{published-rank-anchored} bootstrap used downstream (centre at the published ranks, inject only edge-sampling spread), so that the data ceiling is not mistaken for sampling uncertainty and used to attenuate P1.

\subsection{What the ORCID rebuild certifies: method, not ceiling}

\paragraph{$\rho = 0.74$ against a benchmark of $0.91$ is method certification.} To check that the SpringRank-on-hiring-network \emph{method} (rather than the specific Wapman edges) recovers academic prestige, we rebuild the network from an entirely independent source --- ORCID person-affiliation transitions --- keeping only US doctoral-education$\to$faculty-employment edges, and re-run the identical canonical SpringRank ($\alpha=0.5$, binary weights). On $10{,}687$ academia-level edges over $298$ institutions this ORCID-derived SpringRank reproduces Wapman's published academia ranks at
\[
\mathrm{Spearman}\bigl(s^\star(\text{ORCID edges}),\ \text{rank}^{\text{Wapman}}_{\text{academia}}\bigr) \;=\; 0.74,
\]
while the \emph{public-Wapman benchmark} --- the canonical SpringRank on Wapman's own public edges scored against Wapman's own published academia ranks --- is itself only $0.91$.

\paragraph{Reading.} The correct yardstick for the ORCID rebuild is not $1.0$ but the public-Wapman benchmark, because that benchmark is the most any reconstruction \emph{from public edges} could achieve when reproducing the published academia ordering. An independent network --- different people, different source, different field-tagging --- built entirely outside the Wapman pipeline lands at $0.74$, i.e. within $\approx 0.17$ of a ceiling that is already below $1.0$. That is a certification of the \emph{method}: two unrelated edge sources passed through the same SpringRank optimiser produce nearly the same prestige hierarchy, so the hierarchy is a property of who-hires-whom in US academia, not of any one dataset's idiosyncrasies. Reading $0.74$ as a ``failure to hit $1.0$'' would be a category error --- $1.0$ is unreachable in principle here, since even Wapman's own public edges only hit $0.91$. The certification is strongest exactly where it matters: it holds \emph{with the elite-private top intact} (Stanford, Yale, Harvard, MIT all present as high-prestige employers), which is precisely the part of the hierarchy that the PSEO earnings source truncates (next subsection). Per-field ORCID correlations are lower ($0.45$--$0.68$) and fall \emph{monotonically with edge count} (academia $10.7$k edges $\to 0.74$; per-field $200$--$600$ edges $\to 0.45$--$0.68$), which localises the per-field shortfall to thin per-field edge counts and a crude keyword field-tag --- both data-coverage limits, both fixable --- rather than to the method. The ORCID rebuild is therefore the independent validation that licenses using a SpringRank reconstruction as the AR axis at all.

\subsection{The employer axis and the CIP join}

\paragraph{The join key.} The two axes are linked institution-by-institution and field-by-field through the Classification of Instructional Programs (CIP) code, which is the common vocabulary of fields across the Wapman taxonomy, College Scorecard, and PSEO. A crosswalk maps Wapman's $30$ curated fields to $4$-digit CIP codes and thence to the earnings cells; the gap for a field is then computed only over institutions that appear in \emph{both} the hiring network and the earnings source for that field's CIP. The grain of the join (4-digit CIP) is what makes the suppression and truncation numbers below binding, because disclosure rules and population coverage operate at exactly that cell level.

\paragraph{Scorecard coverage is $\sim$69\% and, crucially, prestige-flat.} College Scorecard covers $\sim$69\% of Wapman institutions, and the decisive fact is that this coverage is \emph{even across prestige quartiles}: from the most to the least prestigious quartile of Wapman institutions the Scorecard match rate is
\[
74.2\% \;/\; 64.5\% \;/\; 71.7\% \;/\; 65.6\% \quad (\text{top}\to\text{bottom}).
\]
\paragraph{Reading.} These four numbers are the match rate of Scorecard against the Wapman institution list within each prestige quartile (Q1 = most prestigious). What matters is not the level ($\sim$69\%) but the \emph{flatness}: there is no monotone gradient, and the top quartile ($74.2\%$) is in fact the best-covered, not the worst. This is what protects the headline gap from being a coverage artifact. If Scorecard systematically dropped, say, low-prestige institutions, then within each field the surviving institutions would be a prestige-truncated sample, the within-field prestige$\leftrightarrow$earnings correlation would be computed on a compressed prestige range, and the gap would be mechanically distorted in a prestige-selected way. Even quartile coverage means the institutions entering each field's Spearman are \emph{not} selected on the very axis (prestige) that defines one side of the gap, so the disagreement we measure is a real prestige$\leftrightarrow$pay disagreement and not an echo of who Scorecard happened to include. This is precisely why Scorecard, not PSEO, is the \emph{primary} earnings source.

\paragraph{PSEO coverage is $\sim$29\% and top-truncated --- so it is a cross-validator, not the primary source.} Census PSEO covers only $\sim$29\% of Wapman institutions and, fatally for a prestige-conditioned statistic, just $21.5\%$ of the top prestige quartile; roughly $18$ of the top-$20$ elite private institutions (Caltech, MIT, Harvard, Princeton, Stanford, Berkeley, Yale, Columbia, Chicago, Penn, Cornell, CMU, $\dots$) are entirely absent.

\paragraph{Reading.} The contrast with Scorecard is the whole point. PSEO's $21.5\%$ top-quartile coverage against Scorecard's $74.2\%$ means PSEO sees the hierarchy with its top sawn off. Because the gap is a within-field correlation \emph{between prestige and pay}, measuring it on a sample that is missing the high-prestige end truncates one of the two ranked variables exactly where its variance lives, which would bias any prestige-conditioned statistic toward whatever pattern survives among the non-elite remainder. PSEO is missing these institutions not at random but \emph{on prestige} (it is a UI-wage, public-institution-skewed source, and the elite privates are disproportionately the missing ones), so a PSEO-based gap would be measured on a structurally different, top-truncated hierarchy. The correct role for PSEO is therefore as an \emph{independent cross-validation and flows} source --- a completely different population definition (UI-employed graduates, real 2023 dollars) against which to check that the Scorecard-based gap \emph{replicates} on the overlapping non-truncated institutions --- not as the primary axis. This is the construction-level reason the cross-source test of \S7 is framed as replication-on-overlap rather than re-measurement.

\paragraph{Why the analysis is bachelor's-level only: $18{,}442$ doctoral cells, $0$ released.} PSEO graduate-level earnings cannot rescue the truncation or extend the gap to higher degrees, because the graduate cells are disclosure-suppressed. At the $4$-digit-CIP grain there are $18{,}442$ doctoral (degree-level $17$) institution$\times$CIP earnings cells and $0$ are released ($0\%$); the master's level is likewise $0$ released.
\paragraph{Reading.} The pair $(18{,}442,\,0)$ is not ``the data does not exist'' --- the cells are enumerated, so the underlying graduates are there --- it is ``$100\%$ of them are withheld by disclosure-avoidance rules'' at the field grain (PSEO releases graduate earnings only at the coarser $2$-digit-CIP level, which cannot separate, e.g., chemistry from physics). The distinction matters because it determines the \emph{reason} the paper is bachelor's-only: it is a censoring constraint, not a gap in the world, and no amount of additional matching or recoding can recover field-level doctoral earnings from a source that suppresses every such cell. This single number is what forecloses the v2 plan of measuring ER ``at the level where each field's talent lands'' (PhD-ER for pipeline fields) and pins the entire analysis to the bachelor's level.

\paragraph{The $\ge 10$-institution threshold and the $28/30$, $29/30$ counts.} A field can only contribute a gap if its within-field Spearman is computed over enough institutions to be non-degenerate. Requiring at least $10$ matched institutions, Scorecard yields $28$ of $30$ usable BA fields and PSEO $29$ of $30$.
\paragraph{Reading.} The $\ge 10$ threshold is the minimum at which a per-field Spearman is meaningful rather than an artifact of a handful of points: with three or four institutions a single noisy earnings figure can flip the sign of the correlation and hence the gap, so a per-field rank statistic on $<10$ institutions is essentially undefined. The threshold buys \emph{non-degeneracy} --- every field that enters the gap map has a Spearman estimated on a sample large enough that no single institution dominates it. That Scorecard clears it for $28/30$ fields (only astronomy and neuroscience, with $3$ and $1$ matched institutions, fall out) and PSEO for $29/30$ shows the join is broadly populated and the field universe is not built from a few well-covered fields. Note these counts measure only \emph{sample adequacy}; they are upstream of the separate \emph{reliability} gate (signal-fraction and bootstrap-CI screen) that later trims the $28$--$30$ adequately-sized fields to the $\sim$14 used for the headline decomposition --- adequacy of sample size and adequacy of signal are two distinct filters, and the $28/30$ and $29/30$ numbers are the first of them.

\paragraph{Closing the loop.} Taken together these numbers fix the construction: AR is a cardinal SpringRank score recomputed from public hiring edges (justified by the blend's need for differences, certified independently at $\rho=0.74$ against a $0.91$ ceiling, and honest about its $\rho\approx0.77$ from-scratch data ceiling); ER is Scorecard bachelor's field-of-study median earnings (primary, because its $\sim$69\% coverage is prestige-flat at $74.2/64.5/71.7/65.6\%$), with PSEO as an independent cross-validator (relegated from primary by its $21.5\%$ top-quartile coverage and forced to the bachelor's level by the $18{,}442$/$0$ doctoral suppression); and the two axes are joined per institution and field on $4$-digit CIP, with each field admitted only where $\ge 10$ institutions match ($28/30$ Scorecard, $29/30$ PSEO). Every later result rests on these construction choices and inherits their certified limits.\section{The gap definition, the diagnostic battery, and the reliability filter}\label{sec:companion-gap}

This section is about \emph{measurement validity}: before any modelling of the gap can be trusted, we must show that $\mathrm{gap}_f$ is a real object and not a manufactured artifact of the way it is constructed. The main paper states the verdict in two sentences; here we set out the estimator, the failure mode it is exposed to, the four-test battery that diagnoses that failure mode, and the per-field filter that the diagnostic forces upon every downstream analysis. The throughline is a single contrast: the filter is exactly what separates a genuine high-gap field (Biology, which survives) from a noise-dominated one (Chemistry, which is screened).

\subsection{The gap definition and why it takes this form}

For a field $f$ observed across institutions $i$, the academic--employer reputation gap is
\[
\mathrm{gap}_f \;=\; 1 - \mathrm{Spearman}_i\!\left(\mathrm{prestige}_{i,f},\ \mathrm{earnings}_{i,f}\right).
\]

\paragraph{Reading.}
The quantity $\mathrm{Spearman}_i(\cdot,\cdot)$ is the rank correlation, \emph{taken over the institutions $i$ that teach field $f$}, between the academic-reputation score $\mathrm{prestige}_{i,f}$ (the continuous SpringRank reconstruction of the faculty-hiring network of \citet{wapman2022quantifying,debacco2018physical}) and the labour-market earnings $\mathrm{earnings}_{i,f}$ (College Scorecard field-of-study median earnings, \citealp{scorecard}). The construction makes three deliberate commitments. \emph{First, why a rank statistic.} A Spearman correlation depends only on the orderings of the two variables, not their levels or units. This buys three things at once: it is monotone-invariant, so any monotone re-scaling of either prestige (which has no natural cardinal unit) or earnings leaves it unchanged; it is insensitive to the enormous level differences in earnings \emph{across} fields (engineering pays far more than education in dollars) because those differences never enter a within-field ranking; and it is robust to the heavy right tail and heteroskedastic scale of an earnings distribution, which would distort a Pearson correlation. What it assumes is that the substantive question is ordinal --- ``do the two valuation systems agree on \emph{which} schools rank where?'' --- rather than cardinal. \emph{Second, why the $1-(\cdot)$ form.} A raw Spearman of $+1$ means the prestige ordering and the pay ordering coincide perfectly (full integration); subtracting it from $1$ flips the scale so that $\mathrm{gap}_f=0$ is perfect prestige--pay agreement and \emph{larger} values mean \emph{more} decoupling. This orients the measure so that ``high gap'' reads directly as ``the two markets disagree,'' which is the object of interest. \emph{Third, why within-field across-institution.} The correlation is computed holding the field fixed and varying the institution. This is what makes the gap a statement about \emph{decoupling within a discipline} --- among the schools that teach $f$, does the academy's ranking of them line up with the labour market's? --- and explicitly \emph{not} a statement about whether one field outpays another. The cross-field pay hierarchy is differenced away by construction; only the within-field rank agreement survives.

\paragraph{The failure mode this construction is exposed to.}
Precisely because the gap is built from a Spearman over institutions, it inherits a known pathology. Where the cross-institution earnings have little \emph{true} dispersion, or where the cohorts behind each institution's median are small, the earnings \emph{ranking} is dominated by sampling noise. A rank correlation computed against a near-noise ordering is mechanically attenuated toward zero, which through the $1-(\cdot)$ form mechanically inflates $\mathrm{gap}_f$ toward one --- a \emph{high gap by construction, even under perfect market agreement}. So a large $\mathrm{gap}_f$ is ambiguous: it can mean genuine prestige--pay disagreement, or it can mean the earnings signal was too weak to rank reliably. The entire diagnostic battery exists to break this ambiguity field by field.

\subsection{The diagnostic verdict and the global-artifact test}

The battery (null-model placebo, bootstrap, residualization, restriction; \texttt{DIAGNOSTIC\_RESULT.md}) returns the verdict \textsc{real signal}, qualified by the requirement of a per-field reliability filter. Two headline numbers carry that verdict.

\paragraph{The null band, and ``$15\%$ inside, $65\%$ above.''}
The null band is a \emph{perfect-agreement-plus-noise} placebo. For each field we synthesise earnings that rank-track prestige \emph{perfectly} (so the true gap is zero by construction) and then corrupt them with realistic per-institution sampling noise calibrated to each cohort's size (median standard error $1.2533\cdot\alpha\cdot\mathrm{earn}/\sqrt{\mathrm{cohort}}$, with the within-school individual-earnings coefficient of variation $\alpha=0.6$; $1000$ draws per field). The $95\%$ spread of the resulting synthetic gaps is the band of gap values a \emph{perfectly integrated} field could exhibit \emph{purely} from sampling noise. Against this band, only \textbf{15\%} of fields' observed gaps fall \emph{inside} it and \textbf{65\%} sit \emph{above} it. The interpretation is direct and decisive: for roughly two-thirds of fields the observed disagreement \emph{exceeds} anything that pure noise on a perfectly-integrated field could manufacture, so that excess is real decoupling, not an artifact. Only the $15\%$ inside the band are candidly indistinguishable from the placebo. The across-field correlation between observed and synthetic gaps is $0.49$ (Pearson) / $0.65$ (Spearman) --- positive, confirming that noise \emph{does} contribute and must be modelled, but far from one, confirming that noise is not the whole story. This finding is robust to the one free knob: sweeping $\alpha\in\{0.4,0.6,0.8\}$ moves the covered fraction only over $15$--$35\%$ and the mean null gap over $0.19$--$0.51$ against an observed mean of $0.48$, so ``most fields exceed noise'' does not hinge on the noise calibration.

\paragraph{The artifact channel $R^2=0.20$, and why a fifth.}
The third test regresses the gap directly on its two mechanical contaminants --- the cross-institution earnings coefficient of variation (\texttt{earn\_cv}, low CV being exactly the low-true-spread condition that inflates the gap) and a grad-school-pull term (signal-carrying graduates leaving for a PhD):
\[
\mathrm{gap}_f \;\sim\; \mathrm{earn\_cv}_f \;+\; \text{grad-school pull}_f, \qquad R^2 = 0.20.
\]

\paragraph{Reading.}
The two regressors are the explicit machinery of the artifact story: if the gap were a global low-SNR artifact, these two variables --- which encode ``how little earnings signal is there'' and ``where did the signal-carrying graduates go'' --- would absorb most of its cross-field variance. A linear regression is the right estimator here precisely because we want the \emph{share of variance attributable} to the artifact channel, which is exactly what the coefficient of determination reports; it assumes only that the contaminant enters approximately linearly, which is adequate for a variance-accounting verdict rather than a structural estimate. The number $R^2=0.20$ is the load-bearing one: the global-artifact story explains only \emph{a fifth} of cross-field gap variance, leaving four-fifths unexplained by the contaminants. Equally important, the regression's \emph{residuals retain structure} --- Earth Sciences $+0.21$, Philosophy $+0.22$, Biology $+0.17$, Chemistry $+0.17$, Physics $+0.15$ remain high after netting the contaminants, while engineering, statistics, economics and computer science stay negative. A residual that is both low-$R^2$ \emph{and} structure-laden is the textbook signature of \emph{real signal}: a pure artifact would leave a high $R^2$ and structureless residuals. (Notably, computer science does not stand out in this residual, $z=-0.67$ --- its low gap is ``as expected'' for a high-CV field --- yet the null-model test independently certifies that CS's gap is genuine disagreement, not the trivial consequence of high dispersion. The two tests are complementary, not redundant.) The verdict \textsc{real signal} is therefore the conjunction of three facts pointing the same way: most observed gaps exceed the noise band, the artifact channel explains only a fifth of the variance, and what it leaves behind is structured.

\subsection{Bootstrap reliability and the two-regime structure that mandates a filter}

The verdict ``real signal in the aggregate'' does \emph{not} license interpreting every individual field's gap. The second test resamples institutions $1000\times$ within each field and reports the bootstrap distribution of the underlying prestige--earnings Spearman (not the gap; a clearly-positive Spearman is the well-measured, integrated case). The pattern across the gap range is the crux of the whole measurement-validity argument.

\paragraph{The well-measured low-gap anchor.}
At the integrated (low-gap) end the correlations are tight and clearly positive: Computer Science $\mathbf{0.71\,[0.62,0.78]}$ and Statistics $\mathbf{0.75\,[0.52,0.87]}$. These are the fields where the cross-institution earnings spread is large enough to rank schools reliably, so the prestige--pay agreement is measured with precision and the confidence interval sits well away from zero. The CS interval is the tighter of the two ($N$ is large and the earnings signal fraction is high, $0.92$ in the null-model table); Statistics is wider ($[0.52,0.87]$) because its institution count is small ($N=28$), which is the generic reason small-$N$ fields carry wide bootstrap intervals regardless of where their point estimate sits.

\paragraph{Why the high-gap end is wide.}
At the decoupled (high-gap) end the intervals blow out: Physics $\mathbf{0.35\,[0.01,0.62]}$ and Chemistry $\mathbf{0.31\,[0.09,0.50]}$, with Earth Sciences ($0.29\,[-0.00,0.57]$) actually \emph{touching zero}. This width is not incidental --- it is the direct fingerprint of the construction-level failure mode. These fields have near-zero true earnings signal (the null-model table assigns Chemistry and Earth Sciences a signal fraction of $0.02$), so the underlying Spearman is estimated against an almost-random earnings ranking; every resample reshuffles a noise-dominated ordering and the bootstrap distribution spreads across most of $[0,1]$. A wide interval that touches or nearly touches zero means we \emph{cannot distinguish that field's prestige--earnings correlation from zero}, and through the $1-(\cdot)$ map we therefore cannot distinguish its high gap from a pure artifact. This is also why the highest-gap-five and lowest-gap-five intervals \emph{overlap}: the extreme gaps can be coarsely separated from the integrated cluster but cannot be finely rank-ordered among themselves.

\paragraph{The decisive contrast: Biology survives, Chemistry is screened.}
The reason a blanket ``high gap = unreliable'' rule would be wrong is Biology: $\mathbf{0.37\,[0.22,0.52]}$ at $\mathbf{N=138}$. This interval is \emph{tight} and lies \emph{entirely below} the CS interval --- a genuinely low, well-measured prestige--earnings correlation, backed by a large institution count and a healthy signal fraction ($0.71$ in the null-model table, observed gap $0.63$ against a null gap of only $0.15$). Biology is thus a \emph{real} high-gap field: its decoupling is measured precisely and is not noise. Chemistry, with an indistinguishable point estimate ($0.31$) but a signal fraction of $0.02$ and an observed gap ($0.69$) that sits \emph{below} its own null gap ($0.88$), is the opposite --- a noise-dominated artifact. Two fields with nearly the same headline gap fall on opposite sides of the validity line. No single threshold on the gap value could tell them apart; only a reliability statistic can.

\paragraph{The reliability filter the diagnostic mandates.}
This two-regime structure is what forces a \emph{per-field reliability gate} ahead of every downstream analysis: a field is admitted only if its \textbf{signal fraction $\ge 0.5$} \emph{and} its \textbf{bootstrap CI excludes the null band}. The two conditions are complementary --- the signal fraction screens fields whose earnings ranking is noise by construction, and the CI condition screens fields whose Spearman cannot be told from the perfect-agreement placebo. Applied to the data, the gate cleanly partitions the high-gap end: Biology (tight CI, healthy signal fraction) passes and is interpreted as genuine decoupling; the \emph{low-earnings-signal cluster} --- Chemistry, Earth Sciences, and the low-CV engineering fields --- fails on both counts and is screened out. (For these very-low-CV fields the $\alpha=0.6$ noise model in the null test actually \emph{over}-states the noise, so they are most honestly labelled \emph{low-reliability} rather than confirmed artifacts; either way they do not clear the gate.) The substantive payoff is twofold. First, the motivating ``chemistry counterexample'' dissolves: chemistry's headline-high gap is among the \emph{least} reliably measured quantities in the dataset (signal fraction $0.02$), so it never was a clean test of anything and must not be read as market segmentation. Second, and decisively for the paper's credibility, the filter is not a device for discarding inconvenient fields --- it admits the high-gap field that is well measured (Biology) and rejects the high-gap field that is not (Chemistry), establishing that ``high gap'' survives as a real, interpretable object precisely where the earnings signal is strong enough to support it. Everything the paper subsequently builds on the gap --- the discipline decomposition, the licensing channel, the cross-source replication --- is run on the screened, reliable field universe, not the raw one.\section{The headline multilevel decomposition}
\label{sec:companion-decomp}

This section opens the box on the paper's central descriptive claim --- that the academic--employer reputation gap is, first and foremost, a property of \emph{which broad discipline area} a field belongs to, not of the field's own idiosyncrasies. The paper (\S4) states the headline as a single line --- $\tau^2=0.0231$, $\mathrm{ICC}=0.30$, $\eta^2=0.50$ --- and a forest of cluster means. Here we explain why each modelling choice was made, what the estimator buys and assumes, and why every number takes the value it does. The companion claims here trace to \texttt{CLUSTER\_DECOMP\_RESULT.md} (script 14) and \texttt{GRANULARITY\_RESULT.md} (script 13).

\subsection{The unit of analysis and why it must be re-weighted}

\paragraph{The input is a vector of noisy field-level gaps.} The decomposition does not see institutions or earnings. Its input is one number per reliable fine field $f$ --- the gap $\mathrm{gap}_f = 1 - \mathrm{Spearman}_i(\text{prestige}_{i,f}, \text{earnings}_{i,f})$ of \S4 --- \emph{together with} a sampling standard error $\mathrm{SE}_f$ for that gap, taken as the bootstrap-CI half-width (floored at $0.05$ so that no single field can claim implausibly infinite precision). The reliability filter of \S4 leaves $n=48$ fine fields entering the headline meta-model. This is the first design fact to internalise: each data point is an \emph{estimate with a known imprecision}, and an honest decomposition of its variance must separate the part that is real dispersion across fields from the part that is merely sampling chatter. A field whose gap is measured on $17$ institutions (Communication Disorders, $\mathrm{SE}=0.233$) cannot be allowed to speak as loudly as Computer Science measured on $161$ ($\mathrm{SE}=0.080$); ignoring that asymmetry would inflate the apparent between-field spread with pure noise.

\subsection{The precision-weighted multilevel meta-model}

\paragraph{The model.} We fit a two-level random-intercept meta-model that regresses each field's gap on a grand mean plus a discipline-cluster random effect:
\[
\mathrm{gap}_f \;=\; \mu \;+\; u_{c(f)} \;+\; \varepsilon_f,
\qquad
u_{c}\sim\mathcal{N}(0,\tau^2),
\qquad
\varepsilon_f\sim\mathcal{N}(0,\mathrm{SE}_f^2),
\]
where $c(f)$ is the CIP-2 (two-digit Classification of Instructional Programs) parent of fine field $f$. In the notation the paper uses, this is the mixed model $\mathrm{gap}_f \sim 1 + (1\,|\,\text{CIP-2})$, fit by restricted maximum likelihood (REML). Stacked over fields, its marginal covariance is
\[
V \;=\; \mathrm{diag}\!\left(\mathrm{SE}_1^2,\dots,\mathrm{SE}_n^2\right) \;+\; \tau^2\,S,
\qquad
S_{fg}=\mathbf{1}\!\left[\,c(f)=c(g)\,\right].
\]

\paragraph{Reading.} The symbols: $\mu$ is the grand-mean gap across all disciplines; $u_{c}$ is the deviation of discipline-cluster $c$ from that grand mean, drawn from a single population of clusters with variance $\tau^2$ (the \emph{between-cluster} variance, the quantity the headline is really about); $\varepsilon_f$ is field $f$'s sampling error around its own cluster's true mean, with \emph{known, field-specific} variance $\mathrm{SE}_f^2$; $S$ is the same-parent indicator matrix, equal to $1$ for any two fields sharing a CIP-2 parent and $0$ otherwise. The matrix $V$ is the load-bearing object. Its diagonal, $\mathrm{diag}(\mathrm{SE}_f^2)$, is the \emph{measurement} (sampling) variance and is what implements precision weighting: GLS/REML under this $V$ multiplies each field's residual by $1/\mathrm{SE}_f^2$, so noisy small fields are automatically down-weighted and cannot dominate the fit. The off-diagonal term, $\tau^2 S$, encodes the \emph{structure}: every pair of fields under the same CIP-2 parent shares a common between-cluster variance component $\tau^2$, which is exactly a \emph{compound-symmetry} (exchangeable) block --- within each discipline block the marginal correlation between any two fields is constant, and across blocks it is zero. That block structure is the formal statement of "fields in the same discipline are correlated because they inherit the same cluster intercept."

Why \emph{this} estimator over the obvious alternatives? A raw one-way ANOVA (or, equivalently, a fixed effect per cluster) treats every cluster mean as estimated independently and on equal footing; with $18$ clusters of which several contain a single fine field (Computer/Info, History, Psychology, English, Architecture each have $n_{\text{fields}}=1$), the unpooled cluster means for those singletons are wildly imprecise, and ANOVA would credit their noise to "between-cluster" variance. A pure pooled regression (one $\mu$, no random effect) goes to the opposite error and assumes the clusters are identical, discarding the very structure we want to measure. The random-intercept model is the principled middle: \emph{partial pooling}. It treats the $18$ cluster means as themselves a sample from a population, estimates how dispersed that population is ($\tau^2$), and shrinks each observed cluster mean toward the grand mean by an amount that grows with that cluster's imprecision. What this \emph{buys}: a between-cluster variance that is purged of sampling noise and of small-cluster over-fitting, plus regularised ("borrowed-strength") cluster means that are usable even for singletons. What it \emph{assumes}: that cluster effects are exchangeable draws from a single Normal population (a regularisation prior, defensible for a descriptive variance budget but not a causal claim), that the supplied $\mathrm{SE}_f$ are approximately correct, and that CIP-2 is the right level of aggregation --- an assumption the paper stress-tests by re-running at the recovered $66$-field universe and by the measurement correction (both in \S7).

\paragraph{The between-cluster variance: $\tau^2 = 0.0231$, SD $0.152$ gap-units.} REML returns $\tau^2=0.0231$, i.e.\ a between-cluster standard deviation of $\sqrt{0.0231}=0.152$ in gap units. Decoding the magnitude: the gap itself ranges over roughly $[0,1]$ in principle and, on the shrunk cluster means, from $0.348$ (Computer/Info) to $0.753$ (Health) --- a spread of about $0.40$. A between-cluster SD of $0.152$ means that moving "one typical discipline" away from the grand-mean discipline shifts the expected gap by about a sixth of the entire observed range. That is a large structural signal: discipline membership alone, before any covariate, moves the gap by a substantial fraction of its full swing. Note that this $\tau^2=0.0231$ is the \emph{covariate-free} headline from script 14; the companion-adjacent model in \texttt{GRANULARITY\_RESULT.md} (script 13), which additionally conditions on dispersion and a licence dummy, reports a closely concordant parent-cluster variance of $\tau^2=0.0248$ (SD $0.158$). The two agree to within rounding precisely \emph{because} the discipline structure is not an artefact of those covariates --- adding dispersion and licence does not absorb the between-cluster variance, which is itself part of the paper's argument that the structure is first-order.

\subsection{The intraclass correlation and why it is not $\eta^2$}

\paragraph{The ICC.} The headline share is the intraclass correlation,
\[
\mathrm{ICC} \;=\; \frac{\tau^2}{\tau^2 + \overline{\sigma^2_{\text{within}}}} \;=\; 0.30,
\]
where $\overline{\sigma^2_{\text{within}}}$ is the mean within-cluster variance of the gap.

\paragraph{Reading.} The symbols: $\tau^2$ is the between-cluster variance just estimated; $\overline{\sigma^2_{\text{within}}}$ is the average dispersion of field gaps \emph{around their own cluster's mean}; their sum is the total gap variance, partitioned into a between part and a within part. The ratio is therefore the \emph{fraction of total cross-field gap variance attributable to differences between broad discipline areas}. The value $\mathrm{ICC}=0.30$ is the paper's headline number: roughly $30\%$ of all the variation in field-level gaps is explained simply by knowing which CIP-2 discipline the field sits in; the remaining $\sim 70\%$ is within-discipline (field idiosyncrasy plus residual sampling noise). This is what licenses the substantive claim "the gap is primarily a discipline-area phenomenon" --- not because $30\%$ is a majority, but because a \emph{single categorical variable with $18$ levels} accounting for nearly a third of the variance dwarfs any of the continuous covariates examined in \S5 (the within-cluster observability gradient, below, is a far weaker second-order effect). The ICC is the precision-weighted, partial-pooled quantity: $\tau^2$ in its numerator has already been shrunk and noise-purged by the meta-model.

\paragraph{The one-way $\eta^2$ and why it differs: $\eta^2 = 0.50$.} As a deliberately unregularised cross-check, the same partition is computed as a raw one-way effect size,
\[
\eta^2 \;=\; \frac{\mathrm{SS}_{\text{between}}}{\mathrm{SS}_{\text{total}}} \;=\; \frac{\sum_{c} n_c\,(\bar{g}_c-\bar{g})^2}{\sum_{f}(g_f-\bar{g})^2} \;=\; 0.50,
\]
the share of the total sum of squares lying between the (\emph{raw}, unshrunk) cluster means.

\paragraph{Reading.} Here $\bar{g}_c$ is the \emph{raw} mean gap of cluster $c$ (no pooling, no precision weights), $n_c$ its number of fields, and $\bar{g}$ the overall mean; $\mathrm{SS}_{\text{between}}$ and $\mathrm{SS}_{\text{total}}$ are the ordinary between-group and total sums of squares of a one-way ANOVA. The functional form is the textbook unconditional effect size: every field contributes equally, every cluster mean is taken at face value, and nothing is shrunk. It is reported precisely \emph{because} it is the naive counterpart to the ICC, and the contrast is the point.

Why do $\mathrm{ICC}=0.30$ and $\eta^2=0.50$ \emph{differ}, and which is right? They differ because they are two different estimators of "between-share," computed under opposite attitudes toward noise. The $\eta^2$ is a raw one-way ratio: it does not down-weight imprecise fields and, crucially, it does not shrink small or noisy clusters toward the grand mean. Singleton and high-SE clusters therefore enter $\eta^2$ at their full, unregularised observed means --- and an unregularised mean of a noisy small cluster sits \emph{farther} from the grand mean than its true value, mechanically inflating $\mathrm{SS}_{\text{between}}$. The ICC, by contrast, is built from a $\tau^2$ that has already pulled those noisy small clusters back toward the grand mean (partial pooling) and weighted every field by $1/\mathrm{SE}_f^2$; that shrinkage removes the sampling-noise component that masquerades as between-cluster spread in $\eta^2$. Hence $\mathrm{ICC} (0.30) < \eta^2 (0.50)$ is exactly what theory predicts: the gap from $0.50$ down to $0.30$ is the noise-and-overfit that regularisation strips out, not a contradiction. The two should be read as a \emph{bracket} --- the conservative, noise-purged share is $0.30$, the raw upper-bound share is $0.50$ --- and the paper takes $0.30$ as the headline precisely because it is the disciplined estimate. (For completeness, and to keep three numbers conceptually distinct: a \emph{third} figure of $0.45$ appears in \S7 as the \emph{measurement-corrected} ICC. That is yet another object --- it corrects the \emph{within}-cluster variance for sampling noise rather than the between, and it moves $0.30\to0.45$ because subtracting the mean $\mathrm{SE}^2$ from the within term shrinks the denominator. It must not be conflated with either $0.30$ (raw partial-pooled ICC) or $0.50$ (raw $\eta^2$); the ordering $0.30<0.45<0.50$ is internally coherent, with $0.45$ sitting between the two precisely because it noise-corrects the within term but, unlike $\eta^2$, retains partial pooling on the between term.)

\paragraph{The cluster counts: $48$ fields, $18$ clusters, $23$ locked-grain units.} Three counts pin down the scope. The meta-model is fit on $n=48$ \emph{reliable fine fields} --- the survivors of the \S4 reliability filter, the only fields whose gaps are well enough measured to enter the variance budget. These distribute across $18$ CIP-2 \emph{clusters} (the random-effect levels: Computer/Info, Math \& Stats, Social Sci, Engineering, $\dots$, Health), which is why $\tau^2$ is estimated from $18$ exchangeable draws --- a modest number that itself motivates partial pooling rather than fixed effects. Separately, the granularity exercise reports $23$ \emph{locked-grain reliable units}: $16$ fine-grain fields that pass reliability on their own \emph{plus} $7$ rolled-up CIP-2 parents formed by pooling thin children that individually failed the filter (e.g.\ the Biological Sciences parent at $149$ institutions, or the Health parent at $104$). The $23$ is the reporting grain for the cluster-mean table and the pooled observability test; the $48$ is the estimation grain for $\tau^2$. They differ because roll-up trades resolution for coverage --- it rescues $28$ previously-dropped fields (Aerospace Engineering, Chemistry, Sociology, Statistics, $\dots$) into the analysis by aggregating them to a reliable parent, buying breadth without manufacturing signal (the result file is explicit that pooling improves "coverage and rigor, not the information").

\subsection{The integrated--decoupled typology}

\paragraph{The two ends, decoded.} The shrunk, partial-pooled CIP-2 means order the disciplines on a single axis from \emph{integrated} (prestige predicts pay, low gap) to \emph{decoupled} (prestige and pay diverge, high gap). At the integrated end: Computer/Info $0.348$, Math \& Stats $0.352$, Social Science $0.421$, Engineering $0.432$. At the decoupled end: Health $0.753$, Arts $0.746$, Natural Resources $0.730$, Biological Sciences $0.678$, Physical Sciences $0.661$. Substantively, the integrated end is the cluster of fields whose academic hierarchy and labour-market pay ordering largely \emph{coincide}: in computing, mathematics/statistics, economics-adjacent social science, and engineering, the institutions the academy ranks highest are also the ones whose graduates the market pays most, so the within-field prestige$\to$pay Spearman is high and the gap ($1-\mathrm{Spearman}$) is low. The decoupled end is where the two valuation systems come apart. Two distinct mechanisms populate it, and the typology is honest that they are not one thing: the \emph{credential areas} (Health most sharply, with the rolled-up Health parent's raw gap reaching $1.058$) are decoupled because occupational \emph{licensing} compresses the pay distribution --- a licensed nurse or audiologist earns the licensed wage almost regardless of the prestige of the granting institution, so prestige cannot transmit to pay and the gap is mechanically high (this is the channel isolated as $b_{\text{licensure}}=+0.66$ in \S5); the \emph{natural-science and arts} clusters (Biological, Physical Sciences, Arts, Natural Resources) are decoupled for a different reason --- bachelor's-terminal earnings in these fields are weakly tied to the research-prestige hierarchy that the faculty-hiring network measures, because that hierarchy prices PhD-pipeline academic standing rather than the labour-market value of a terminal bachelor's degree. The key methodological point the companion must stress is that this entire typology is \emph{read off the cluster structure}, not hand-coded: the licence lesson and the PhD-pipeline lesson \emph{emerge} as the high-gap end of an unsupervised CIP-2 decomposition, which is far stronger evidence than if the high-gap fields had been labelled by hand.

\paragraph{Shrinkage made concrete: Biology raw $0.830$ vs.\ shrunk $0.678$.} The single cleanest illustration of partial pooling in the table is the Biological Sciences cluster, whose \emph{raw} mean gap is $0.830$ but whose \emph{shrunk} (partial-pooled) mean is $0.678$ --- a pull of $0.152$ toward the grand mean. Decoding: the raw $0.830$ is the unpooled average of a cluster of five fine fields, some of them noisily measured; the model judges that part of that high value is sampling noise and small-sample over-extension, and regularises it back toward the overall mean by an amount calibrated to the cluster's imprecision. The shrunk $0.678$ is the \emph{regularised} estimate the paper reports --- the model's best guess at the cluster's \emph{true} mean having borrowed strength from the global population of clusters --- while $0.830$ is the unpooled, credulous estimate. The direction is diagnostic: extreme raw means always shrink \emph{inward}, and the more imprecise the cluster, the harder the pull (the result files note per-parent shrinkage factors ranging $0.42$--$0.94$, with high-SE fields borrowing most). This is also why the shrunk decoupled-end values ($0.678$, $0.661$, $0.730$) are systematically \emph{less} extreme than their raw counterparts ($0.830$, $0.686$, $0.960$): shrinkage compresses the tails of the cluster-mean distribution, which is precisely the same operation that drives $\mathrm{ICC}=0.30$ below $\eta^2=0.50$. The two facts are the same fact viewed at the cluster level versus the variance level.

\subsection{The within-cluster observability gradient: a real but second-order modifier}

\paragraph{The secondary question.} Having established that most structure is \emph{between} clusters, the natural follow-up is whether anything systematic happens \emph{within} a cluster: does a field that is more "observable" --- operationalised as lower earnings dispersion across institutions, the idea being that when pay is tightly clustered employers can read productivity more directly and lean less on the prestige credential --- exhibit a lower gap? Two estimates address this.

\paragraph{The pooled and within-cluster estimates.} The pooled (inclusive) rank association across all $23$ reliable units is
\[
\mathrm{Spearman}\!\left(\mathrm{gap},\ \text{dispersion}\right) \;=\; -0.44 \quad [-0.74,\ +0.01],
\]
and, controlling out the cluster means, the within-cluster regression slope is
\[
\frac{\partial\,\mathrm{gap}_f}{\partial\,\text{dispersion}_f}\;=\;-0.527 \quad [-1.08,\ +0.02] \qquad (n=46).
\]

\paragraph{Reading.} The first statistic is a rank correlation between a field's gap and its cross-institution earnings dispersion, pooled over all reliable units \emph{without} removing cluster identity; the second is the partial slope of gap on dispersion in a model with CIP-2 fixed effects, so it asks the cleaner within-discipline question --- "among fields in the \emph{same} discipline, does the more observable one have a lower gap?" --- on the $46$ fine fields that admit it. A Spearman is used for the pooled view because the relationship is monotone-but-not-necessarily-linear and robust to the long right tail of dispersion; an OLS slope with cluster dummies is used for the within view because we want a magnitude net of the dominant between-cluster structure, not merely a rank. Both point the \emph{same} direction --- negative --- which is the theoretically predicted sign: more observable fields (lower dispersion) have \emph{smaller} gaps, consistent with employers relying less on the prestige signal when productivity is easier to read. The slope magnitude $-0.527$ says that, within a discipline, a one-unit increase in dispersion is associated with roughly a half-unit increase in the gap.

But the decisive feature is the confidence interval: \emph{both} CIs span (or all but touch) zero --- $[-0.74,+0.01]$ for the pooled Spearman and $[-1.08,+0.02]$ for the within slope. That is why the paper demotes this to a "weak second-order modifier, not the headline." The point estimates are real and correctly signed and the central tendency is unlikely to be a fluke given two independent operationalisations agree, but the data cannot exclude a true effect of zero at conventional confidence, so it would be overclaiming to present observability as an \emph{established} driver alongside the between-cluster decomposition. The contrast in evidentiary status is the lesson: the between-cluster structure ($\mathrm{ICC}=0.30$, $\tau^2$ SD $0.152$) is a first-order, robust finding; the within-cluster observability gradient is a directionally consistent but CI-fragile refinement on top of it. The companion model in \texttt{GRANULARITY\_RESULT.md} reinforces this --- under full precision weighting the dispersion slope attenuates to $-0.044\,[-0.08,-0.00]$, "only marginally bounded," and the parallel licence shift collapses from a large unweighted OLS value to $+0.063\,[-0.05,+0.18]$ once the two highest-SE licence fields (Communication Disorders, Nursing) are properly down-weighted and low-SE Accounting pulls the licence mean back. The honest reading is that neither covariate creates the structure; the discipline cluster does, and the observability gradient is a faint, correctly-signed second-order overlay that the paper reports plainly rather than promotes.\section{Channel decomposition: the one clean channel and the redundant rest}\label{sec:companion-channels}

The main paper states that ``\emph{one} clean external channel --- occupational licensing ($b\approx+0.66$) --- plus a large irreducible residual'' explains the gap, and that ``additional candidate channels are collinear and add essentially nothing.'' This section explains how that single sentence is earned: what was regressed on what, why those two anchors and not others, why the second anchor's null is informative rather than an artifact, and why a battery of three further channels --- each individually plausible, two individually significant --- evaporates the moment they are asked to work together. The throughline is that the hunt for mechanism, conducted honestly, returns redundancy: licensing is the only external lever that moves the gap and survives every control, and everything else is either the same lever in disguise or orthogonal noise.

\subsection{The precision-weighted channel regression}

\paragraph{The estimating equation.} We project the field-level gap onto two externally-measured, field-aggregate anchors built from the American Community Survey (ACS PUMS), estimated by precision-weighted least squares:
\[
\mathrm{gap}_f \;=\; \beta_0 \;+\; b_{\text{licensure}}\,\ell_f \;+\; b_{\text{absorption}}\,a_f \;+\; \varepsilon_f,
\qquad
\widehat{b} \;=\; \big(X^\top W X\big)^{-1} X^\top W\,\mathrm{gap},
\quad
W = \mathrm{diag}\!\big(1/\mathrm{SE}_f^2\big).
\]

\paragraph{Reading.} The outcome $\mathrm{gap}_f = 1 - \mathrm{Spearman}_i(\text{prestige}_{i,f},\text{earnings}_{i,f})$ is the within-field prestige-pay rank disagreement defined in the main text. The two regressors are the \emph{external anchors}. \emph{Licensure intensity} $\ell_f$ is the precision-weighted share of field $f$'s employed degree-holders who work in occupations requiring a licence (the STRICT construction: near-universal-licensure occupations only --- health practitioners SOCP 29-1xxx, lawyers 23-1xxx, K--12 teachers 25-2xxx, architects, psychologists), a number in $[0,1]$. \emph{Academic absorption} $a_f$ is the share of field $f$'s \emph{doctorate}-holders employed as postsecondary teachers (SOCP 25-1xxx), also in $[0,1]$. Because both anchors run from 0 to 1, each coefficient $b$ reads directly as ``the change in the gap as the channel goes from 0\% to 100\% of the field.'' The estimator is WLS with weights $W=\mathrm{diag}(1/\mathrm{SE}_f^2)$: each field's gap is itself a Spearman estimated on a finite, sometimes small, set of institutions, so its sampling variance differs wildly across fields (Computer Science is measured on $N{=}138$ institutions; some fine fields on a dozen). Ordinary least squares would let a noisily-measured small field swing the slope as hard as a precisely-measured large one; precision weighting down-weights the noisy fields by exactly their imprecision, which is the same logic that produces the REML between-cluster variance in the discipline decomposition of the main paper --- we re-use it here so the channel regression and the variance decomposition treat measurement error identically.

Why \emph{these two} anchors, and only two? They operationalise the two most mechanical, theory-named compression stories. Licensing is the canonical wage-compression mechanism: where a credential is legally required to practise (nursing, law, teaching, pharmacy), the licence floor and ceiling squeeze the pay distribution, so prestige cannot map onto pay even when the academic hierarchy is sharp --- this should \emph{raise} the gap. Academic absorption operationalises the rival ``pipeline'' story: a field whose doctorates mostly stay inside the academy (the natural sciences) might have a prestige hierarchy built for an internal labour market that the external earnings data never sees --- this too could mechanically inflate the gap. The design isolates these two because each is measurable \emph{outside} the prestige-earnings correlation itself, on a public universe, so neither is a transform of the gap. What the regression \emph{buys} is a clean linear attribution: the fraction of cross-field gap variance that co-moves with two named, external mechanisms. What it \emph{assumes} --- and we state this as forcefully as the result file does --- is exactly nothing causal. The anchors are field aggregates joined to a field-level outcome, so every association is \emph{ecological} (field-level, never within-field); the anchors are themselves consequences of the same discipline structure that drives the gap (a collider/over-control risk); and the ACS undergraduate-field universe on which $\ell_f, a_f$ are measured is not the Wapman PhD-field or Scorecard BA-earnings universe of the gap. This is a \emph{decomposition} --- a linear projection / variance attribution --- not an IV estimate: the anchors are not instruments and no exclusion restriction is claimed.

\paragraph{Reading.} On the $n{=}47$ fields with both anchors (one of the 48 non-degenerate gap fields lacks an absorption anchor), the primary specification (strict licensure, ACS absorption) returns
\[
b_{\text{licensure}} = +0.66\ [+0.33,\,+0.97], \qquad
b_{\text{absorption}} = -0.02\ [-0.42,\,+0.46], \qquad
R^2 = 0.28, \qquad
\mathrm{corr}(\ell,a) = -0.21.
\]
The licensing coefficient $+0.66$ means that moving a field from 0\% to 100\% licensed occupations \emph{raises} its gap by 0.66 gap-units --- an enormous effect against a gap scale that runs roughly 0 to 1.3. The sign is the predicted one (regulation compresses pay, decoupling it from prestige) and the confidence interval $[+0.33,+0.97]$ excludes zero cleanly, so licensing is the one channel that survives as a CI-excludes-0 mechanical lever. The value is large because the channel is doing real work in the most decoupled corner of the data: Health ($\ell=0.58$) and Education ($\ell=0.61$) are the two highest-licensure clusters, and they are precisely the clusters whose gap collapses most under netting (below). $R^2=0.28$ says that just over a quarter of cross-field gap variance is linearly attributable to these two mechanical channels jointly; the remaining $\sim$72\% is the irreducible decoupling residual that the main paper bills as ``large.'' That residual is the headline finding, not a failure: two clean external mechanisms cannot explain most of the gap, which is what makes the gap a substantive object rather than a mechanical accounting identity.

\paragraph{Reading.} The absorption coefficient $-0.02$ with CI $[-0.42,+0.46]$ is a genuine null: the point estimate is essentially zero and the interval is symmetric about zero. The non-obvious move is to certify that this null is \emph{informative} rather than an artifact of the absorption signal being absorbed by a correlated licensing term. That certification is the number $\mathrm{corr}(\ell,a)=-0.21$. In a regression, one regressor's coefficient can be driven to zero when it shares variance with another (its unique contribution is small only because the partner has already claimed the common variance); a null on a regressor that is strongly collinear with a significant one tells you little. Here the two anchors are \emph{near-orthogonal} ($|{-}0.21|$ is weak), so licensing has \emph{not} stolen absorption's explanatory variance --- there is almost no shared variance to steal. The pipeline channel, as measured, simply does not predict the gap. This reading is reinforced by the sensitivity grid: under the SDR doctorate-only absorption proxy the coefficient even flips to $-0.35$ (on a tiny $n{=}17$ SEH-only sample), which is the opposite sign --- a channel that cannot hold its sign across two reasonable measurements is not a stable mechanism. So the absorption null is a real, reportable negative result: of the two named mechanical stories, only licensing is a lever. We report the strict/broad licensure split openly ($b_{\text{licensure}}=+0.66$ strict vs $+0.45$ broad, with engineering the swing field once Professional-Engineer licensure is counted) rather than selecting the specification that flatters the headline.

\subsection{The redundant rest: three channels that add essentially nothing}

\paragraph{Setup.} Having found one clean lever, the honest next question is whether \emph{more} external channels add explanatory power. We test two further field attributes against the gap, each conceptually distinct from it and measured on ACS BA-holders: an industry-versus-academic \emph{pay-wedge} (median private-for-profit minus academic/education-sector earnings) and a \emph{public-sector share} (fraction of the field's employment in government). The joint standardized model is
\[
\mathrm{gap}_f = \beta_0 + b_w\, z(\text{pay-wedge}_f) + b_p\, z(\text{public-sector}_f) + b_\ell\, z(\text{licensure}_f) + \varepsilon_f,
\]
with $z(\cdot)$ the standardization that puts the coefficients on a common scale, again precision-weighted, on $n{=}48$ fields.

\paragraph{Reading.} Standardizing each channel to unit SD makes the three slopes directly comparable in ``gap-units per channel-SD,'' which is the right scale for asking \emph{which} channel dominates when all compete. The functional form is again a linear precision-weighted projection rather than anything causal, for the same ecological reason as before --- but here the explicit purpose is a horse-race, not an attribution, so the quantity to read is the \emph{joint} $R^2$ and the \emph{combined}-compression story, not the individual slopes. The result file is emphatic on this point and we follow it: the three channels share a mechanism (compression), so their individual coefficients are not separately identified well; reporting any one slope in isolation would over-read it.

\paragraph{Reading.} The multi-channel model returns
\[
R^2_{\text{multi}} = 0.29 \quad\text{vs}\quad R^2_{\text{licensure-only}} = 0.27,
\]
with standardized slopes
\[
b_\ell = +0.11\ [+0.07,\,+0.14], \qquad b_w = -0.03, \qquad b_p = -0.01.
\]
The entire story is in the gap between $0.27$ and $0.29$: adding \emph{two} new channels to the licensing-only baseline buys two percentage points of $R^2$. That is the operational meaning of ``add essentially nothing.'' Inside the joint model, only licensing survives standardization --- $b_\ell=+0.11$ with a tight CI $[+0.07,+0.14]$ that excludes zero --- while the pay-wedge ($-0.03$) and public-sector ($-0.01$) coefficients have collapsed to the neighbourhood of zero. They collapse not because they are individually meaningless (they are not --- see the bivariate decode below) but because they are \emph{collinear compression mechanisms}: they encode the same underlying regulatory/sector-pay compression that licensing already captures, so once licensing is in the model they have no unique variance left to explain. This is the textbook signature of redundancy: a bivariate signal that vanishes under joint control because the controlled-for variable was carrying it all along.

\paragraph{Reading.} Two further numbers confirm that the added channels are inert. First, the cluster gap ranking computed on the multi-channel residual is \emph{identical} to the one computed on the licensing-only residual: $\mathrm{Spearman}=+1.00$. A Spearman of exactly $+1.00$ between two residual rankings is the strongest possible statement of redundancy --- it means that after netting licensing, adding the pay-wedge and public-sector channels does not move a single cluster relative to any other. The decoupled end (Natural Resources, Arts, Physical Sciences) stays decoupled; the integrated end (Math \& Stats, Computer/Info, History) stays integrated; the residual ordering is untouched. Second, the residual standard deviation barely moves: $0.228 \to 0.225$ (from a raw $0.250$, so licensing alone did the work of shrinking $0.250\to0.228$, and the two extra channels shave only a further $0.003$). The residual SD is the dispersion of the part of the gap that the channels \emph{cannot} explain; that it scarcely shrinks means the large structured residual that licensing leaves behind is essentially untouched by the additional channels. Together the $+1.00$ Spearman and the near-static SD say: the added channels are not a new mechanism, they are a collinear re-description of the one we already had.

\paragraph{Reading.} The cleanest illustration of redundancy-as-disguise is the public-sector channel. Bivariately it is \emph{significant and positive}:
\[
\mathrm{Spearman}(\mathrm{gap},\text{public-sector}) = +0.37\ (p=0.009,\ n=48).
\]
Read alone, this looks like a real second mechanism: government pay scales compress earnings, so a field with more public-sector employment should have a higher gap. But in the joint model its coefficient is $-0.01$ --- it collapses to zero once licensing is netted. The reconciliation is the channel correlation structure: public-sector share correlates $+0.40$ with licensure (and $-0.54$ with the pay-wedge), so the public-sector ``signal'' was \emph{licensing in disguise} --- the licensed fields (teaching, health) are exactly the public-sector-heavy ones, and once licensing is in the model there is nothing left for public-sector to explain. A significant bivariate that vanishes under control is precisely how a confounded-with-licensing channel should behave; reporting only the $+0.37$ would have manufactured a phantom second mechanism. The channel correlations run up to $|0.54|$ across the compression channels, which is why the result file insists on reading the joint $R^2$ rather than the separate slopes.

\paragraph{Reading.} The pay-wedge is the most interesting of the failed channels because, unlike public-sector, it is \emph{not} merely licensing in disguise. Bivariately,
\[
\mathrm{Spearman}(\mathrm{gap},\text{pay-wedge}) = -0.33\ (p=0.022,\ n=48),
\]
and crucially it \emph{survives a field-size control}: partialling out log institutions-per-field moves the relative-wedge correlation from $-0.27$ to $-0.29$ --- if anything slightly \emph{stronger}, not weaker. This is the diagnostic that distinguishes a real-but-redundant channel from a spurious artifact. The negative sign is mechanistically sensible: a larger industry-over-academic pay premium signals sharper market demand, which sharpens market valuation, which makes pay track prestige more closely, which \emph{lowers} the gap. So the pay-wedge is a genuine, externally-measured, field-size-robust correlate of the gap. Why, then, does it not count as a second channel? Because its bivariate $-0.33$ attenuates to a joint standardized $-0.03$ for a reason the field-size check rules out: the attenuation is \emph{collinearity with licensing} ($\mathrm{corr}=-0.28$ with licensure, $-0.54$ with public-sector), not a size confound. It is a real but weaker compression story that licensing already substantially captures; on the proper BA construction it is $-0.33$, roughly half the magnitude of the $-0.61$-class doctorate-only SDR wedge that an earlier, small-sample SEH-only cross-level construction had inflated. The honest verdict is therefore ``holds bivariately, attenuated, redundant jointly'' --- a real second correlate that nonetheless adds two points of $R^2$ and zero re-ordering.

\paragraph{The honest omission: unionization.} One compression channel that \emph{should} be in this battery is unionization --- collective bargaining is a canonical wage-compression mechanism that would sit alongside licensing and public-sector. It is absent because no BLS/CPS union-coverage-by-occupation table exists in the repository and the analysis round forbade importing new collaborator data, so the channel is left \emph{unbuilt rather than fabricated}. We flag it as a documented data gap (it would enter via a crosswalked field$\to$SOC union-coverage file) rather than silently dropping it or inventing a coefficient; given that the two compression channels we \emph{do} measure both collapse into licensing, the prior is that a union channel would behave the same way, but we do not assert a number we did not compute.

\paragraph{What the section establishes.} Read end to end, the channel decomposition delivers an asymmetric result that is more credible for being asymmetric. \emph{One} external anchor --- licensing --- is a robust lever: $b_{\text{licensure}}=+0.66$ in the two-anchor decomposition, a tight standardized $+0.11$ in the four-way horse-race, sign-stable, CI-excludes-0, and the sole channel that re-orders the decoupling ranking. Every other candidate fails for a \emph{diagnosable, reported} reason: academic absorption is a true null on near-orthogonal regressors ($\mathrm{corr}=-0.21$, so the null is informative); public-sector is a significant bivariate ($+0.37$) that is licensing in disguise (collapses to $-0.01$); the pay-wedge is a real, field-size-robust bivariate ($-0.33\to-0.29$) that is nonetheless collinear-redundant jointly ($-0.03$); and unionization is an honest data gap. The joint $R^2=0.29$ and the residual ranking's invariance (Spearman $+1.00$, SD $0.228\to0.225$) are the quantitative seal: beyond licensing, the search for mechanism finds collinear redundancy, not new explanatory power --- which is exactly why the gap's large residual is the paper's substantive object rather than an artifact waiting for the right covariate.\section{Three ruled-out alternatives (scope conditions)}\label{sec:companion-ruleout}

The main paper reports, in \S5.2, three deflationary analyses that test intuitively attractive explanations of the gap and rule each of them out. This section explains the design of those three tests, why each was built the way it was, and --- decoding every number in the ledger --- why each verdict came out as it did. The framing is deliberate: a finding that an explanation \emph{fails} is informative only if the test was capable of detecting it had it been true, and only if the failure is interpretable as a \emph{scope condition} (a statement about what our data \emph{can} and \emph{cannot} resolve) rather than a null of unknown origin. All three analyses inherit the same structural limitation, which is the recurring villain of this section: College Scorecard reports \emph{medians} of observed workers, and the elite-tail outcomes where brand effects are theorised to live \citep{chetty2023diversifying} are invisible to a median. Every named program in these analyses is illustrative only; the robust object is always a distribution or a cross-field statistic.

\subsection{(a) The gap is not generic-brand versus field-specific prestige}

\paragraph{The alternative being tested.} The most natural competing story for a measured prestige--pay wedge is a \emph{mismatched-signal} story: students chase the institution's overall ``brand'' (the academia-wide reputation $G$), while pay actually tracks \emph{field-specific} academic strength (the field reputation $F$) --- the ``Georgia-Tech-CS beats Yale-CS'' intuition. If true, the gap would be an artifact of using the wrong prestige signal, and switching from $G$ to $F$ would close it. The test must therefore put the two signals into a horse race for the same outcome (median field earnings $y$) and ask which one pay actually tracks.

\paragraph{The horse-race quantities.} For each field $f$ we form two within-field rank correlations of prestige with pay across institutions,
\[
c_F \;=\; \mathrm{Spearman}_i\!\left(F_{i,f},\, y_{i,f}\right), \qquad
c_G \;=\; \mathrm{Spearman}_i\!\left(G_{i,f},\, y_{i,f}\right),
\]
and define field reputation's predictive advantage as
\[
\mathrm{ADV}_f \;=\; c_F - c_G .
\]

\paragraph{Reading.} $c_F$ is the slope of pay on \emph{field-specific} prestige: how well the field's own academic ranking of institutions orders their graduates' pay. $c_G$ is the slope of pay on the \emph{generic} academia-wide brand: how well the institution's overall reputation orders pay \emph{within the same field}. $\mathrm{ADV}_f$ is the head-to-head margin: positive when the field signal beats the brand, negative when the brand beats the field signal. The test logic is direct --- if pay tracked field-specific prestige (the alternative), $c_F$ should systematically beat $c_G$ and $\mathrm{ADV}$ should be positive. We use Spearman rather than a level regression here because the gap itself is a rank statistic (\S2), so the horse race must be on the same rank footing to be commensurable; this buys insensitivity to field-level pay scale and to monotone transformations of either prestige axis, at the cost of discarding cardinal premium information (recovered separately in the $\beta$ regression below). One symbol is load-bearingly \emph{not} new: by construction $c_F = 1 - \mathrm{gap}_f$ exactly (verified to machine precision, max absolute deviation $10^{-16}$ in \texttt{FIELD\_VS\_GENERIC\_RESULT.md}), so $c_F$ merely restates the gap. The genuinely new, non-tautological object is $c_G$ --- the gap says nothing about whether the \emph{generic} brand predicts field pay --- and hence the contrast $\mathrm{ADV}$. The assumption this buys its leverage from is that field earnings $y$ are a fair common target for both signals; what it cannot do is see the tail, which is exactly where the alternative might survive.

\paragraph{The verdict, decoded.} Across the 57 fields, the mean field slope is $c_F = +0.40$ and the mean brand slope is $c_G = +0.45$, so $\overline{\mathrm{ADV}} = -0.05$: \emph{on medians the generic brand predicts pay at least as well as --- usually slightly better than --- field-specific prestige}. The mismatched-signal alternative is rejected, and the sign of the rejection is informative. $\mathrm{ADV} < 0$ in $42/57$ fields, i.e.\ in roughly three-quarters of fields the brand wins outright; this is not a knife-edge average masking heterogeneity but a broad tilt. The few fields where field reputation \emph{does} lead are applied/vocational (Pharmacy $\mathrm{ADV}=+0.31$, Special Education $+0.22$, Social Work $+0.15$, engineering subfields) --- exactly the fields where a specific licensed or occupational competence, not a halo, is what employers buy --- and pointedly \emph{not} the fields the intuition invokes (in Computer Science $\mathrm{ADV}=-0.04$).

\paragraph{Why the brand wins: the two signals are nearly the same signal.} The reason $c_G$ keeps pace with $c_F$ is that the two prestige axes are not independent contenders. Define
\[
\rho_{PP}[f] \;=\; \mathrm{Spearman}_i\!\left(F_{i,f},\, G_{i,f}\right),
\]
the within-field agreement between field-specific and academia-wide prestige.

\paragraph{Reading.} $\rho_{PP}$ is a \emph{prestige-only} statistic --- it never touches earnings --- so it diagnoses, prior to any horse race, how much room there even is for the two signals to disagree about pay. Its mean is $+0.79$ (median $+0.82$): field and brand prestige are strongly aligned. This decodes the headline. When two predictors are correlated at $+0.79$, they cannot order a third variable very differently, so $c_F$ and $c_G$ are nearly bound to be close, and which one wins on average is decided by a thin residual margin (here, narrowly, the brand). $\rho_{PP}$ also tells us \emph{where} the horse race is even meaningful: it is near-degenerate in Mathematics ($+0.90$), Physics ($+0.94$), Economics ($+0.90$), but genuinely loose in Pharmacy ($+0.12$), Forestry ($+0.31$), Microbiology ($+0.50$) --- and it is precisely in the low-$\rho_{PP}$ fields that $\mathrm{ADV}$ swings hardest (Pharmacy's $\rho_{PP}=+0.12$ co-occurs with its field-reputation advantage of $+0.31$). So the alternative is not merely rejected on average; it is rejected \emph{because the premise that there are two cleanly separable signals to choose between is itself mostly false}.

\paragraph{The named-example check and its honest reframe.} To make the test concrete we ran a data-driven divergence test in Computer Science: hold the field, take the institutions with the most extreme $F_{\mathrm{pct}} - G_{\mathrm{pct}}$ (field-strong/brand-medium vs.\ brand-strong/field-medium), and compare their median pay. The brand-strong schools out-earn the field-strong ones, $\$119{,}643$ vs.\ $\$98{,}970$ (a $\$20{,}674$ median gap in the brand's favour). The number decodes the same way as the cross-field average --- the brand, not the field signal, is where the pay is --- but its \emph{provenance} carries a disclosed caveat that strengthens rather than weakens the conclusion. Schools strong in \emph{both} CS and overall brand (MIT, Stanford, CMU, Berkeley) are excluded by construction, because they are not divergence cases. What remains in the ``field-strong'' arm is not a set of CS-elite-but-brand-weak schools (those barely exist --- CS-elite almost always implies brand-strong) but regional publics with only \emph{median} CS prestige and very low brand (Louisiana--Lafayette, UT El Paso, Mississippi State). The test therefore actually contrasts ``mid-field/low-brand public'' against ``low-field/high-brand private,'' which conflates geography and selection --- and even so the brand arm wins. The contrast is not universal: in Mechanical Engineering the field-strong arm narrowly wins ($\$91{,}385$ vs.\ $\$90{,}088$), consistent with engineering being one of the applied fields where field competence is priced. We report whichever way each table fell; that is the point.

\paragraph{The scope condition.} The right reading is not ``the brand causally pays more.'' It is that \emph{on medians} we cannot distinguish a brand effect from selection (generic prestige $\approx$ selectivity $\approx$ student composition), and that the place where the field-versus-brand distinction would actually bite --- the elite tail (top-1\%, elite firms and graduate schools) --- is by construction absent from median earnings \citep{chetty2023diversifying}. A null (indeed reversed) brand advantage on medians is fully consistent with a large brand payoff concentrated in a tail we do not observe. So the alternative is ruled out \emph{as an explanation of the gap we measure}, with the explicit scope condition that the gap is a median-earnings phenomenon and brand effects live in a tail outside our window. The supporting cardinal evidence (the $\beta$-regression of $y$ on standardized $F$ and $G$) classifies $15$ fields as brand-dominant against only $7$ field-dominant, corroborating the rank-based verdict in dollar terms while remaining subject to the same tail blindness.

\subsection{(b) Integration varies in strength, not in ``shape''}

\paragraph{The alternative being tested.} A second natural story is that the gap, being a single correlation, hides qualitatively different \emph{mechanisms}: some fields might reward prestige \emph{gradually} (a graded ladder), others only at the apex (a \emph{threshold}/winner-take-all premium where only the most prestigious tier earns more), others not at all (flat). If integration came in distinct shapes, the one-dimensional gap would be an over-compression and the interesting structure would be the shape, not the strength. The test must classify each field's prestige$\to$pay relationship into graded/threshold/flat and then ask whether that classification carries information \emph{beyond} the gap itself.

\paragraph{The construction.} For each of the 44 fields with $\ge 20$ institutions we separate the apex from the body: the top-prestige-decile pay premium versus the prestige--pay correlation $c_{\mathrm{rest}}$ \emph{among the non-top-decile} institutions. A field is \textsc{threshold} if the relationship is driven only by the top decile (large apex premium, weak gradient below), \textsc{graded} if it persists in the body ($c_{\mathrm{rest}}$ strong), and \textsc{flat} if neither. The raw classification is $\textsc{graded}\;30 \,/\, \textsc{threshold}\;5 \,/\, \textsc{flat}\;9$.

\paragraph{Reading.} The split is built so that ``shape'' is operationalised by \emph{where in the prestige distribution} the pay response lives, which is the only shape distinction medians can even attempt: $c_{\mathrm{rest}}$ measures the gradient in the body, and the top-decile premium isolates the apex. This buys a concrete, computable taxonomy; what it assumes is that a winner-take-all premium would show up in the \emph{median} of the top decile --- an assumption we flag as false below. The counts are reported but immediately suspect: a sweep over the (flat-cut $\times$ $c_{\mathrm{rest}}$-cut) grid moves \textsc{flat} from $6$ to $13$ and \textsc{threshold} from $1$ to $8$, so the three-way split is cutoff-dependent, not a clean partition.

\paragraph{Why the shapes collapse into strength: the redundancy correlations.} The decisive diagnostic is whether the shape classification is orthogonal to the gap or merely a re-cut of it. Two Spearman correlations settle it:
\[
\mathrm{Spearman}\!\left(c_{\mathrm{rest}},\, c_F\right) = +0.96, \qquad
\mathrm{Spearman}\!\left(\text{class ordinal},\, c_F\right) = +0.80 .
\]

\paragraph{Reading.} The first correlation asks whether the body-only gradient $c_{\mathrm{rest}}$ tells us anything the overall integration strength $c_F = 1-\mathrm{gap}$ does not; at $+0.96$ the answer is essentially no --- $c_{\mathrm{rest}}$ and $c_F$ rank the fields almost identically, so ``how steep is the gradient below the apex'' is the same information as ``how integrated is the field.'' The second correlation asks the same of the categorical label itself: ordering the classes (flat$<$threshold$<$graded, or by integration) tracks $c_F$ at $+0.80$, meaning the \textsc{graded}-versus-\textsc{flat} distinction \emph{is mostly restating whether integration is strong or weak}. A Spearman this close to $+1$ between a derived axis and the original is the formal signature of \emph{redundancy}: the added channel is not a new dimension but a noisy copy of the existing one. This is why the alternative is ruled out --- not because shape does not exist in principle, but because our shape variable does not separate from strength, so we resolve \emph{strength, not shape}. The only cell with genuine claim to being shape-based is \textsc{threshold} (a large apex premium with a weak body gradient is not implied by a given $c_F$), but it is the smallest and most cutoff-sensitive cell ($5/44$, swinging to $8$ under a different cut).

\paragraph{Why threshold is rare: a limitation, not a finding.} The temptation is to read ``only $5$ threshold fields'' as evidence that winner-take-all is rare in higher education. We explicitly do not. A winner-take-all premium, by definition, is an \emph{elite-tail} phenomenon (top-1\%, elite firms), and Scorecard reports the \emph{median} of the top decile, not its tail \citep{chetty2023diversifying}. The median of an apex tier whose payoff is concentrated in its own upper tail will look unremarkable. So median data \emph{cannot test} winner-take-all; the rarity of \textsc{threshold} is what the data cannot resolve, not a substantive null. (An earlier draft mis-framed this rarity as a defence of the graded view; the corrected reading is that it is a ceiling on what medians show.) The scope condition is therefore: integration varies mostly in strength ($\approx$ the gap), with a small, cutoff-sensitive set of apex-only fields layered on top that we are not powered to confirm; we can resolve \emph{how much} prestige is priced, not the \emph{shape} of how it is priced.

\subsection{(c) Program ``mispricing'' is geography and selection, not value-added}

\paragraph{The alternative being tested.} The third story is the most seductive because it promises a policy product: residualize earnings on prestige within a field, and the institutions with large positive residuals are ``underrated'' programs that the academy under-prices but the market rewards --- a value-added discovery engine. The test must surface those residuals and then ask, adversarially, what they are actually made of. We define the standardized within-field residual $r$ so that $r>0$ marks a program paid above its academic standing (market-overperforming) and $r<0$ a program whose academic standing exceeds its market pay.

\paragraph{The diagnostic: regressing the residual on what it might really be.} The key quantity is a variance decomposition of the residual against two competing explanations --- destination geography and academic brand --- run at the institution level:
\[
R^2\big(\,\bar r_{\text{inst}} \sim \text{destination-state FE}\,\big) = 0.46,
\qquad
R^2\big(\,\bar r_{\text{inst}} \sim \text{brand}\,\big) = 0.00 .
\]

\paragraph{Reading.} $\bar r_{\text{inst}}$ is an institution's mean standardized residual across its fields (computed only where $\ge 3$ fields are available, so single-field institutions cannot masquerade as systematic). The first $R^2$ is the share of cross-institution variation in ``mispricing'' explained by a full set of \emph{destination-state} fixed effects --- one dummy per state graduates land in (PSEO flows, matched for $38\%$ of institutions, $44\%$ of programs). The second is the share explained by the academic-hiring brand $G$. The state-FE specification was chosen because geography is the most mechanical confound of an earnings residual: a program is ``market-overperforming'' partly because its graduates work in high-wage labour markets, which has nothing to do with the program's value-added. Fixed effects are the right estimator here because we want to absorb \emph{all} between-state wage-level differences nonparametrically without imposing a functional form on the geography--wage relationship; what this buys is a clean upper bound on geography's contribution, and what it assumes is that the matched PSEO subset is representative enough for the bound to transfer. The contrast of the two $R^2$ values is the entire finding. A state-FE $R^2$ of $0.46$ against a brand $R^2$ of $0.00$ says that \emph{where graduates work explains nearly half of the ``mispricing,'' and academic brand explains none of it}. The residual that the value-added story wants to read as hidden program quality is, to first order, a map of regional wage levels: the named over-performers are selective coastal/urban privates in high-wage states (Santa Clara in Silicon Valley; Georgetown, Dartmouth, Vanderbilt, Boston College) and the named under-performers are publics in lower-wage interior states. The disagreement between the two valuations is mostly \emph{where you are and who you enrol}, not value-added.

\paragraph{Why the residual is barely a residual.} A second number completes the demolition. The standardized residual is $+0.88$-correlated with the \emph{raw} within-field earnings ranking. This follows mechanically from how little prestige explains within-field pay: the mean within-field $\mathrm{Spearman}(F,y)\approx 0.42$, so prestige accounts for only $\sim 18\%$ ($\approx 0.42^2$) of within-field earnings rank-variance. When the regressor you residualize on explains only $\sim 18\%$ of the outcome, the residual is overwhelmingly just the outcome itself, lightly adjusted. Hence a ``market-overperforming'' program is, $+0.88$ of the way, simply \emph{a high-paying program for its field} --- which loops straight back to the geography and selection signal the state-FE decomposition already isolated. The residual does not surface a new latent quality; it re-surfaces pay.

\paragraph{Horizon instability seals it.} If the residuals identified durable program quality they should be stable across earnings horizons. They are not: $29\%$ of program residuals \emph{flip the sign} of their over-/under-valuation between the 1-year and 4-year horizons (rank-stability vs.\ the 4-year measure is only $+0.56$ at 1yr and $+0.64$ at 5yr, i.e.\ $\sim 31$--$41\%$ shared rank-variance). A program label that says ``overrated'' at one horizon and ``underrated'' at another cannot be measuring a fixed value-added; it is measuring noise plus a slowly-moving geography/selection signal. The residual \emph{distribution} is reasonably stable, but \emph{individual program names are not} --- which is exactly why every named program in this analysis is illustrative only and the robust object is the institution-level aggregate and the distribution, never a single program.

\paragraph{The scope condition.} Putting the three numbers together --- geography explains $0.46$ and brand $0.00$; the residual is $+0.88$ pay; $29\%$ of labels flip --- the verdict is that ``where two valuations disagree'' is \emph{not} ``where a program adds value.'' The disagreement is dominated by geography and selection and is horizon-unstable. We make no causal or ``underrated'' claim; ``over-valued''/``over-performing'' name a descriptive gap between two orderings, not a quality judgment. The honest scope condition is that direct selectivity controls (admit rate, SAT/ACT) are absent from the available data, so we cannot net selection fully --- but the PSEO-state check shows the selection confound is empirically load-bearing rather than merely asserted, since geography alone already absorbs far more residual variance than brand. The would-be value-added engine is ruled out as a value-added engine; what survives is a descriptive disagreement map dominated by the very confounds a causal version would have to remove.\section{The two-signal model: built, calibrated, and honestly not validated}\label{sec:model-companion}

The main paper devotes a single guarded section to a generative model and an entire co-author decision note (\texttt{COAUTHOR\_NOTES.md}) to the question of whether it should survive to submission at all. This companion section explains why that guardedness is the correct posture. The model is not a result; it is an \emph{organizing lens}. Its purpose is to show that the empirical regularities documented in \S\S4--5 can all be produced by one small, transparent data-generating process in the Altonji--Pierret employer-learning tradition \citep{altonji2001employer}, and --- this is the harder and more honest half --- to identify the exactly one external prediction by which the model could be \emph{killed}, run that test, and report that it did not pass. We first present and decode the generative equation, then the gap it implies, then the licensing-netted residual that constitutes the model's load-bearing external claim. We then decode, number by number, the calibration, the predictions it does and does not reproduce, the make-or-break validation that fails, and the ICC recalibration that is fitting rather than prediction.

\subsection{The generative earnings equation}

Each field $f$ contains institutions $i$. Each institution's graduates carry two latent qualities, an academic value $\theta_A$ and a market value $\theta_M$, drawn jointly Gaussian with a field-specific alignment $\rho_f$:
\[
(\theta_{A,f,i},\ \theta_{M,f,i}) \sim \mathcal{N}\!\left(0,\ \begin{bmatrix} 1 & \rho_f \\ \rho_f & 1 \end{bmatrix}\right),
\qquad
\rho_f = \mathrm{corr}_i(\theta_A,\theta_M).
\]
Academic prestige observes $\theta_A$ with noise; earnings at horizon $t$ are a credential-weighted blend of the two values that drifts toward market value as employers learn, then compressed by licensing and perturbed by a horizon-decaying sort noise:
\[
\mathrm{prestige}_{f,i} = \theta_{A,f,i} + \sigma_A\,\varepsilon,
\qquad
w(\tau_f,t) = e^{-\tau_f t},
\]
\[
E_{f,i}(t) = (1-\lambda_f)\big[\,w(\tau_f,t)\,\theta_{A,f,i} + (1-w(\tau_f,t))\,\theta_{M,f,i}\,\big] \;+\; \sigma_E(t)\cdot \mathrm{noise},
\qquad
\sigma_E^2(t) = \sigma_\infty^2 + \sigma_0^2\,e^{-\delta t}.
\]

\paragraph{Reading.} Every symbol carries a substantive load. $\theta_A$ is the latent \emph{academic value} of an institution's graduates in field $f$ --- the thing the faculty-hiring network of \citet{wapman2022quantifying} prices; $\theta_M$ is their latent \emph{market value} --- the thing employers ultimately pay for. The single field primitive $\rho_f=\mathrm{corr}(\theta_A,\theta_M)$ is the \emph{field alignment}: the cross-institution correlation between the two valuation systems, so that $(1-\rho_f)$ is the genuine \emph{valuation divergence} --- exactly the wedge $\mathrm{gap}_f$ is meant to detect. $\mathrm{prestige}$ is $\theta_A$ read through observation noise $\sigma_A$, encoding that academia hires on academic value but imperfectly. The credential weight $w(\tau_f,t)=e^{-\tau_f t}$ is the heart of the employer-learning story: at $t=0$ earnings track the credential ($\theta_A$) fully; as the horizon $t$ grows and as the field's \emph{observability} $\tau_f$ is higher, $w$ decays toward zero and earnings migrate from $\theta_A$ onto $\theta_M$, because employers have learned true productivity and no longer need to lean on the credential signal. This is precisely the Altonji--Pierret mechanism \citep{altonji2001employer} in functional form: the weight on the credential decays with observability and tenure. The licensing factor $(1-\lambda_f)$ is a \emph{multiplicative compression}: where graduates funnel into licensed occupations, the cross-institution signal in earnings is squeezed toward a regulated common value, so $\lambda_f\to 1$ erases between-institution earnings dispersion entirely. Finally $\sigma_E^2(t)=\sigma_\infty^2+\sigma_0^2 e^{-\delta t}$ is the sort-noise schedule: early cohorts are unsorted (large $\sigma_0^2$ component), and as $t$ grows the transient noise decays at rate $\delta$ toward an irreducible floor $\sigma_\infty^2$.

Why this form and not another? The exponential weight $w=e^{-\tau t}$ is the minimal monotone kernel that (i) is bounded in $[0,1]$ so it reads as a mixing weight, (ii) collapses two distinct economic forces --- field observability $\tau_f$ and tenure $t$ --- into a single product $\tau_f t$, which is what lets the model speak to both the cross-field observability gradient and the multi-horizon ICC decline with no extra parameters, and (iii) reproduces the Altonji--Pierret prediction that the credential premium decays toward zero rather than reversing sign. The Gaussian latent pair is chosen so that the whole model has a closed-form Pearson proxy (below) that can be checked against the Spearman Monte Carlo, giving an analytic handle on \emph{why} each prediction holds rather than a black-box simulation. What the form \emph{buys}: a three-primitive description $(\rho_f,\tau_f,\lambda_f)$ of any field plus one common noise schedule, from which the gap is a deterministic function --- so cross-field gap variation becomes structure, not idiosyncrasy, which is exactly the discipline-structured ICC finding the model is meant to rationalize. What it \emph{assumes}, and these are real costs: linearity of earnings in the two latent values (no complementarity, no convexity in the credential premium); Gaussian, homoskedastic-across-institutions signals (so the model cannot see the elite tail where \citet{chetty2023diversifying} locate brand payoffs); and licensing entering as a pure multiplicative scalar $(1-\lambda_f)$ rather than, say, shifting the mean or truncating the distribution. The model takes $(\rho_f,\tau_f,\lambda_f)$ as primitives and does \emph{not} explain why they vary by discipline; it is therefore \emph{consistent with} the between-discipline ICC structure of \S4, not a derivation of it.

\subsection{From structural parameters to the observable gap}

The gap is defined on simulated data exactly as in the empirical pipeline, as one minus the within-field rank agreement of prestige and earnings:
\[
\mathrm{gap}_f(t) = 1 - \mathrm{Spearman}_i\!\left(\mathrm{prestige}_{f,i},\ E_{f,i}(t)\right).
\]
The closed-form Pearson proxy, validated against the Spearman Monte Carlo, makes the dependence on the primitives explicit:
\[
\mathrm{corr}(\mathrm{prestige},E) = \frac{(1-\lambda_f)\big[w + (1-w)\rho_f\big]}{\sqrt{(1+\sigma_A^2)\,\big[(1-\lambda_f)^2 B + \sigma_E^2(t)\big]}},
\qquad B = w^2 + (1-w)^2 + 2w(1-w)\rho_f.
\]

\paragraph{Reading.} The left-hand object is the within-field rank co-movement of the two valuation systems; $\mathrm{gap}_f=1-\mathrm{corr}$ so the gap rises as that co-movement falls. The Spearman (rank) form is chosen over a Pearson gap for the same reason as in the empirical paper: it is invariant to the earnings \emph{level} differences across fields (engineering pays more than education) and to any monotone re-scaling, isolating purely whether the prestige ordering and the pay ordering coincide among institutions teaching $f$. The closed form decodes \emph{why} the structural parameters move the observable gap. Three facts fall straight out. First, a \emph{floor}: as $t\to\infty$, $\sigma_E^2\to\sigma_\infty^2$ and the gap settles at a value that increases in $(1-\rho_f)$ and, at $\rho_f=1$, is independent of $w$ and hence of $\tau_f$ --- divergence sets the floor, observability only modulates the approach. Second, a \emph{licensing}$\times$\emph{noise} interaction: if $\sigma_E=0$ the factor $(1-\lambda_f)$ cancels top and bottom because rank correlation is scale-invariant, so licensing moves the gap \emph{only} in the presence of noise --- a non-obvious testable implication, that the licensing slope should decay with horizon as noise resolves. Third, $\tau_f$ acts purely through $w$, with $\partial\,\mathrm{gap}/\partial\tau$ scaling with $(1-\rho_f)$ and vanishing at $\rho_f=1$. What this buys is an interpretive bridge: every empirical regularity is mapped to a sign condition on a partial derivative of one expression. What it assumes is that the Gaussian/Pearson closed form approximates the Spearman gap closely enough to reason with --- the simulation uses Spearman throughout and the two are reported to agree closely, so the closed form is used for intuition, not for the headline numbers.

\subsection{One calibration, no per-fact tuning}

\paragraph{The single parameter vector.} The entire prediction battery is run at \emph{one} held-fixed calibration: $N=150$ institutions per field, $\sigma_A=0.50$, $\sigma_\infty=0.30$, $\sigma_0=1.00$, $\delta=0.70$, with the three field primitives $(\rho_f,\tau_f,\lambda_f)$ swept across fields. This is the single most important methodological fact about the model and it is easy to miss. There is \emph{no per-fact tuning}: the same five global constants generate the licensing prediction P2, the absorption null P3, and the lead--lag direction P5. The discipline of one calibration is what converts the exercise from curve-fitting into a (weak) test --- a model with one free parameter per target can hit any target, whereas a single vector that simultaneously lands several disjoint empirical facts is doing real work. The price of this discipline is visible in the misses below: because nothing is tuned to the ICC-vs-horizon fact, the model gets that fact wrong, and the paper reports it wrong rather than re-tuning.

\subsection{Predictions that match (and the null that matches)}

\paragraph{P2 --- licensing raises the gap, $\approx+0.64$ against the empirical $+0.66$.} The model's licensing slope $\partial\,\mathrm{gap}/\partial\lambda$ is positive at every horizon and declines in $t$ ($+0.58$ at $t{=}1$, $+0.48$ at $t{=}4$, $+0.33$ at $t{=}8$), exactly the licensing$\times$noise interaction the closed form predicts. Crucially the gap--$\lambda$ curve is \emph{convex}, so a linear slope over the observed support $\lambda\in[0,0.85]$ understates the effect; the model's full-support effect at the four-year horizon is $\mathrm{gap}(\lambda{=}0)=0.33 \to \mathrm{gap}(\lambda{=}1)\approx 0.97$, a rise of $\approx +0.64$. \emph{Decoding why this number takes this value:} it is not fit to anything --- it is set entirely by the calibrated noise level $\sigma_E(t)$, since with no noise the licensing factor would cancel out of the rank correlation. That this un-tuned full-support rise ($+0.64$) lands essentially on top of the empirical $b_{\text{licensure}}=+0.66$ from \S5 is the model's cleanest quantitative hit, and the near-coincidence is the reason the licensing channel is the one piece of the model the paper is willing to call durable. The honest caveat travels with it: the quoted slope depends on the $\lambda$-range, which is why the \emph{full-support rise} is reported rather than the smaller local slope.

\paragraph{P3 --- the absorption null, $-0.04$ ($p=0.13$), is the expected outcome.} Academic absorption (the share of graduates entering academia) appears in \emph{none} of the model's primitives $(\rho,\tau,\lambda,\sigma)$. The simplest gap-generating model therefore contains no absorption term, so the empirical null on absorption is the \emph{predicted} result, not an awkward silence. When absorption is drawn independent of the primitives, its simulated gap coefficient is $-0.04\,[-0.09,+0.01]$, $p=0.13$ --- a clean null alongside a live licensing coefficient $b_\lambda=+0.34$. \emph{Decoding why it is null:} the coefficient is statistically indistinguishable from zero precisely because the regressor is, by construction, orthogonal to everything that generates the gap; a non-null here would have \emph{required} an absorption channel and thereby falsified the minimal model. This is what the main text calls ``the cleanest hit'': the model both reproduces the empirical null and converts it into a falsifiable condition --- the empirical null implies absorption $\perp (1-\rho)$. The contrast confirms the regressor is not inert by construction: when absorption is instead generated \emph{confounded} with divergence, it acquires a spurious $+0.26$ ($p\approx 0$). The null is therefore real and informative, not a dead variable.

\paragraph{P5 --- gap-change correlates $+0.71$ with divergence.} In the two-wave simulation ($t{=}1$ vs $t{=}5$), prestige is the \emph{slow prior} (the same latent $\theta_A$ across waves) while earnings \emph{drift} toward $\theta_M$ as $w(\tau,t)$ decays. The change in the gap correlates $+0.71$ with the divergence $(1-\rho_f)$: low-divergence fields' gaps \emph{close} (mean $\Delta=-0.10$) while high-divergence fields' gaps \emph{open} (mean $\Delta=+0.09$). \emph{Decoding the $+0.71$:} it is large and positive because divergence is the only force that pushes the drifting earnings series \emph{away} from the stationary prestige series; where $\theta_A$ and $\theta_M$ are aligned the drift moves earnings along prestige and the gap cannot open. This is the model's directional prediction for the eventual lead--lag test of \S8: earnings should \emph{lag} prestige, and the decoupling should be largest where $(1-\rho)$ and $\tau$ are largest. Together with the licensing channel, this lead--lag direction is the model's second --- and only other --- piece of durable, testable content.

\paragraph{P1 and P4 --- partial.} The horizon prediction P1 (gap falls toward a divergence-set floor) matches in aggregate --- floors rise with divergence ($(1-\rho)=0.1\to$ floor $0.24$; $0.4\to0.49$; $0.7\to0.74$) and the panel mean falls $0.50\to0.42$ --- but the \emph{most divergent tail} is a disclosed partial miss: in the default calibration, w-drift can beat noise-resolution there, so those gaps \emph{rise} $1\to5$ yr, whereas the data show even high-gap clusters falling. P4 (observability gradient) reproduces the mechanism and the strong $\tau\times(1-\rho)$ interaction but \emph{not} a weak $\tau$ \emph{main} effect at 4--5 yr in the default calibration. These two are reported as partial, not sold as matches, which is the honest reading: three clean reproductions (P2, P3, P5), two partial.

\subsection{The make-or-break external test, and why it leaves the model unvalidated}

The model's one load-bearing external prediction is an estimator claim. Netting the licensing channel (and, per the model's own recommendation, \emph{not} netting the non-mechanism absorption term) should leave a residual that estimates valuation divergence:
\[
r_f \;=\; \mathrm{gap}_f - b_{\text{licensure}}\cdot\text{licensure}_f
\;=\; \mathrm{gap}_f - (+0.650)\cdot\text{licensure}_f
\;\;\widehat{=}\;\; \big(1-w(\tau_f)\big)\,(1-\rho_f),
\]
so that $r_f$ should correlate \emph{positively} with any independent proxy for the divergence $(1-\rho_f)$.

\paragraph{Reading.} The left side subtracts the empirically estimated licensing contribution from the gap; the model's closed form says what is left is the floor term $(1-w(\tau_f))(1-\rho_f)$ --- genuine valuation divergence scaled by how far earnings have drifted off the credential anchor. This is the move that promotes the residual from ``whatever is orthogonal to two anchors'' (a descriptive, un-identified object in \S5) to an \emph{estimand}: an estimate of $(1-\rho_f)$. The estimator is chosen --- net $\lambda$, do not net absorption --- because the model proves netting $\lambda$ is unbiased \emph{iff} $\lambda \perp \rho$, and that netting absorption can only cost a degree of freedom or, worse, bias the residual if absorption happens to proxy $\rho$ (over-control on a non-mechanism). What the claim buys is enormous if true: it gives the irreducible residual of \S5 a structural meaning. What it assumes is the orthogonality condition $\mathrm{corr}(\text{licensure},\text{divergence})\approx 0$, which the model insists is empirically checkable rather than asserted --- and which is exactly what the validation tests. \emph{What would confirm it:} a positive correlation between $r_f$ and an external academic-versus-market divergence proxy. \emph{What would refute it:} a null or, worse, contradictory signs across independent proxies.

\paragraph{B1 is UNCONFIRMED, and this is the honest headline.} Two independent divergence proxies were constructed: an O*NET skill-content proxy (applied minus abstract skill intensity of a field's occupations) and an SDR salary-wedge proxy ($|\text{academic}-\text{industry salary}|/\text{mean}$). The model predicts both should correlate \emph{positively} with the residual. They do not agree: O*NET gives $\mathrm{Spearman}(r,\text{proxy})=+0.41\,[+0.08,+0.68]$ ($n=48$) while SDR gives $-0.61\,[-0.88,-0.11]$ ($n=17$). \emph{Decoding why this combination is fatal to the ``estimated'' reading:} the two proxies are not two noisy reads of one construct --- they are themselves \emph{negatively correlated}, $-0.64$, so they measure different things and disagree in sign. One is positive (confirming), one is robustly negative (refuting), and on the same 17 fields where both exist O*NET is $+0.75$ while SDR is $-0.61$. There is no principled, non-circular basis to prefer one over the other. The B3 orthogonality check, intended as a tie-breaker, does not rescue O*NET: $\mathrm{corr}(\text{licensure},\text{O*NET})=-0.03$ (clean, licenses the netting for that proxy) but $\mathrm{corr}(\text{licensure},\text{SDR})=+0.39$ (materially non-zero) --- and yet partialling licensure out of the SDR relation leaves it at $-0.57$, so SDR's contradiction is independent of any licensure contamination. \emph{Decoding the directional disagreement substantively:} a residual that is the \emph{same} object cannot correlate $+0.41$ with one valid divergence measure and $-0.61$ with another; the sign split means at least one proxy does not measure the cross-institution academic-versus-market value correlation $\rho_f$ the model means (O*NET measures a task-abstractness \emph{level}, SDR a doctorate-level price wedge cross-level to the bachelor's gap), or the residual is not that divergence. Either way the make-or-break external prediction is \emph{not borne out}. Two further adversarial caveats deepen the downgrade: netting licensing is nearly inert here (the residual $\approx$ the raw gap at Spearman $+0.87/+0.90$, so B1 is effectively testing the raw gap, not a licensing-purged quantity), and the O*NET $+0.41$ partly reflects a field-size confound (it attenuates to $+0.34$ controlling institutions-per-field). The verdict is therefore that the residual-as-divergence reading is \emph{downgraded from ``estimated'' to ``unconfirmed,''} and the model is \emph{not validated}. The paper does not paper over this with P2/P3/P5: those are internal-consistency reproductions, whereas B1 is the one external bet, and it failed.

\subsection{The ICC recalibration is fitting, not prediction}

The model's most explicit \emph{internal} miss is the ICC-versus-horizon trajectory. Empirically the between-cluster ICC falls with horizon ($0.50$ at 1 yr to $0.30$--$0.33$ at 4--5 yr; \S4,\,\S7). The default calibration gets the \emph{grand} gap falling but its between-field variance slightly \emph{rises} ($0.0172\to0.0200$ across horizon) --- a direct contradiction, reported as an honest miss rather than hidden. A disclosed recalibration makes the early sort-noise $\sigma_0$ discipline-varying and tied to licensing (licensed fields compressed hardest at 1 yr, resolving by 5 yr). Under that recalibration the between-field variance \emph{falls} ($0.0297$ at 1 yr $\to 0.0246$ at 5 yr), matching the empirical decline, and it does so \emph{without} breaking the three clean predictions: P2 holds (the $t{=}4$ licensing slope is $+0.604$, up from the default $+0.318$), P3 holds (absorption $b_{\text{absorption}}=+0.029$, $p=0.19$, still null), and P5 holds ($\mathrm{corr}(\Delta\mathrm{gap},1-\rho)=+0.40$, down from $+0.48$ but still positive).

\paragraph{Reading.} The mechanism is economically plausible: licensed/heavily-compressed disciplines (Health, credentialed fields) have very large but \emph{structured} 1-yr gaps near $1.0$ that resolve as cohorts sort by 5 yr, so between-field spread is largest early --- which is what the empirical 1-yr picture shows (Arts $\approx 1.3$ vs CS $\approx 0.3$ already at 1 yr). And the recalibration is honest in two senses: it costs exactly \emph{one} added degree of freedom (a field-specific $\sigma_0$ in place of a common $\sigma_0$), disclosed rather than tuned silently, and it does not damage P2/P3/P5. But the decisive interpretive point is that this is \emph{fitting, not prediction}. The default model, with its single un-tuned calibration, gets the between-variance direction \emph{wrong} ($0.0172\to0.0200$, rising); the only way to reverse it to falling ($0.0297\to0.0246$) is to add the very degree of freedom the default lacked and tie it to the target. That the same parameters happen not to break the other predictions shows the model's \emph{structure} is rich enough to \emph{accommodate} the ICC decline, not that it \emph{predicts} it from first principles. The mechanism is asserted (licensed disciplines hugely compressed at 1 yr) rather than independently measured. Accordingly the recalibration is \emph{not} folded into the headline calibration; it is reported only to establish that the ICC miss is a calibration limitation of one fixed vector, not a structural impossibility of the model.

\subsection{What the model is for}

Read as a whole, the model's ledger is: three clean internal reproductions (P2 licensing, P3 absorption null, P5 lead--lag direction), two partial (P1 horizon tail, P4 observability main effect), one disclosed internal miss that only a fitted extra parameter repairs (ICC-vs-horizon), and one external make-or-break prediction that comes back \emph{unconfirmed} because two independent divergence proxies contradict each other in sign. This is deliberately not enough to call the model validated, and the paper does not. The two pieces that \emph{survive} as durable, testable content are the licensing channel --- whose un-tuned full-support effect ($\approx+0.64$) coincides with the empirical $b_{\text{licensure}}=+0.66$ --- and the lead--lag direction P5 predicts for the dynamic test of \S8. Everything else is an organizing lens: a way to see that one small employer-learning process \citep{altonji2001employer} can generate the discipline-structured wedge, while being explicit that the lens has not earned the status of a result \coauthor{whether to keep \S6, cut to appendix, or invest in validation --- see COAUTHOR\_NOTES.md}.\section{The robustness battery}\label{sec:robust-companion}

The headline of the paper---that the academic--employer reputation gap is structured by broad discipline area, with roughly 30\% of cross-field gap variance sitting \emph{between} CIP-2 clusters---rests on a chain of estimation choices, each of which a hostile referee can attack. This section documents the battery of robustness checks that pre-empt those attacks. We organise it by \emph{threat to validity}: for each check we state what could go wrong, what statistic certifies that it did not, and---the part the main paper omits---\emph{why} that statistic takes the value it does. Two of these checks do more than survive: the measurement-error correction and the crosswalk audit both \emph{strengthen} the headline, and we say so explicitly. None of the checks was outcome-shopped; each is run as a single pre-specified script (\texttt{22\_hardening}, \texttt{24\_crosswalk\_audit}, \texttt{32\_pseo\_gap\_crossval}, \texttt{21\_multihorizon}, \texttt{15--19} for OpenAlex) and reported whole.

\subsection{Statistical hardening of the cluster decomposition}

\paragraph{Threat 1: the integrated--decoupled ranking is an artefact of how many institutions populate each cluster.}
A Spearman correlation computed on few institutions is high-variance, so a cluster's gap could rank high \emph{merely} because it is measured on few institutions---the ranking would then be an artefact of sample size, not of any real prestige--pay structure. The diagnostic is whether the cluster order survives equalising the per-cluster institution count. We pool each CIP-2 cluster (canonical SpringRank on the pooled child-field edges plus cohort-weighted earnings), then subsample \emph{every} eligible cluster down to a common $n=20$ institutions and recompute each cluster gap.

\[
\rho_{\text{rank}}\bigl(\text{gap}^{\,\text{full-}n},\ \text{gap}^{\,n=20}\bigr) = +0.98 \quad(\text{18 clusters}), \qquad
\mathrm{corr}(\text{gap},\,n) = -0.54 .
\]

\paragraph{Reading.} The first quantity is the Spearman rank correlation between each cluster's gap measured on its full institution count (which ranges 23--155, median 100) and the same cluster's gap after every cluster is forcibly down-sampled to exactly $n=20$ institutions; it is taken over the 18 clusters that have at least 20 institutions. The second, $\mathrm{corr}(\text{gap},n)=-0.54$, is the cross-cluster correlation between a cluster's gap and its institution count. We chose forced equal-$n$ subsampling rather than, say, a parametric size control because it directly neutralises the mechanism of concern---Spearman's sampling variance is a function of $n$, so equalising $n$ equalises that variance across clusters, leaving only the structural signal. What it \emph{buys} is a clean falsification target: if the ranking were an $n$ artefact, equalising $n$ would scramble it. What it \emph{assumes} is that pooling a cluster's heterogeneous child fields into one SpringRank yields a meaningful ``cluster gap'' (the cluster gap is itself a construct), and that the institutions retained after subsampling are representative of the cluster.

The decode is the interplay of the two numbers. The $-0.54$ is the referee's ammunition: gap and $n$ \emph{are} negatively correlated, so the integrated (low-gap) end---Computing $n=96$, Math \& Stats $n=105$, Social Sci $n=150$, Engineering $n=126$---does sit among the better-populated clusters, while several decoupled (high-gap) clusters (Architecture $n=29$, Philosophy/Religion $n=31$, Education $n=23$) are thin. A naive reader could conclude the whole ranking is just ``thin clusters look decoupled.'' The $+0.98$ defeats exactly that conclusion: when every cluster is judged on the \emph{same} 20 institutions, the order is essentially unchanged (Health $1.058\to1.040$, Social Sci $0.305\to0.331$ at the two extremes; the within-row shifts are all small and order-preserving). So the negative gap--$n$ correlation is real but is not the \emph{cause} of the ordering---it is at most a partial confound that the equal-$n$ recomputation strips out. The honest residual caveat, carried from the source, is that equal-$n$ controls Spearman's \emph{variance} but not the deeper question of \emph{which} institutions populate each cluster; the $+0.98$ certifies robustness to the first, not the second.

\paragraph{Threat 2: the ICC is contaminated by per-field sampling noise, so ``30\% between disciplines'' is mismeasured.}
The ICC reported in the paper is $\mathrm{ICC}=\tau^2/(\tau^2+\overline{\text{within}})$, the share of cross-field gap variance attributable to between-CIP-2 differences. The within-cluster variance $\overline{\text{within}}$ is estimated from the observed dispersion of field gaps inside each cluster---but each field's gap is itself measured with sampling error (a Spearman on a finite, noisy set of institutions). That error inflates the \emph{observed} within-cluster scatter, biasing the ICC \emph{downward}. The correction removes the sampling component:

\[
\widehat{\text{within}}_{\text{true}} \;=\; \widehat{\text{within}}_{\text{obs}} \;-\; \overline{\mathrm{SE}^2},
\qquad
\mathrm{ICC}_{\text{corr}} \;=\; \frac{\tau^2}{\tau^2 + \widehat{\text{within}}_{\text{true}}} .
\]

\paragraph{Reading.} Here $\widehat{\text{within}}_{\text{obs}}$ is the raw, observed within-cluster variance of field gaps; $\overline{\mathrm{SE}^2}$ is the mean of the per-field gap-bootstrap sampling variances (the average squared standard error of a single field's gap estimate); $\widehat{\text{within}}_{\text{true}}$ is the resulting estimate of the \emph{true} within-cluster dispersion of gaps; and $\tau^2$ is the between-cluster variance from the precision-weighted REML meta-model, which is \emph{already} noise-free because the meta-model's working covariance $V=\mathrm{diag}(\mathrm{SE}^2)+\tau^2\cdot(\text{same-parent})$ explicitly partitions sampling variance out of $\tau^2$. The subtraction is the textbook variance-decomposition identity for an additive measurement-error model: an observed value is true value plus independent sampling noise, so observed variance equals true variance plus mean sampling variance, and the unbiased estimator of true variance subtracts $\overline{\mathrm{SE}^2}$. We chose this subtractive correction over inflating $\tau^2$ or over a fully Bayesian errors-in-variables model because it is transparent and uses quantities we already estimate (the bootstrap SEs); what it \emph{buys} is an ICC that is not biased down by per-field noise, and what it \emph{assumes} is that the bootstrap $\mathrm{SE}^2$ is the \emph{only} noise and that $\tau^2$ contains no model misspecification---if REML's $\tau^2$ has absorbed misspecification, the corrected ICC could be biased \emph{up} (this is the disclosed counter-objection).

\paragraph{Reading.} The estimated quantities substitute as follows.
\[
\widehat{\text{within}}_{\text{obs}}=0.0549,\quad \overline{\mathrm{SE}^2}=0.0267,\quad \widehat{\text{within}}_{\text{true}}=0.0282,\quad \tau^2=0.0231,
\]
\[
\mathrm{ICC}_{\text{raw}}=\frac{0.0231}{0.0231+0.0549}=0.30 \;\longrightarrow\; \mathrm{ICC}_{\text{corr}}=\frac{0.0231}{0.0231+0.0282}=0.45 .
\]
The observed within-cluster variance $0.0549$ is almost half sampling noise: the mean field-gap sampling variance is $0.0267$, so $0.0549-0.0267=0.0282$ is the genuine within-cluster heterogeneity. That nearly halves the denominator's within-term, and \emph{because the denominator shrinks while the numerator $\tau^2=0.0231$ is held fixed}, the ratio rises from $0.30$ to $0.45$. Substantively: roughly half of the apparent ``field-idiosyncratic'' scatter inside a discipline was never real---it was the noise floor of measuring a gap on a finite institution set. Once that floor is removed, $0.45$ of the \emph{true} cross-field gap variance is between disciplines, not $0.30$. This is the sense in which the correction \emph{strengthens} the headline: the raw ICC was conservative, and noise-correction moves the between-discipline share up, not down. The number that does \emph{not} move is $\tau^2$ itself, precisely because the REML weighting had already purged sampling noise from the between-component---so the correction is applied only where noise actually lives, in the within-component.

\paragraph{Threat 3: prestige ranks are treated as fixed, so the gap CIs are too narrow.}
The gap bootstrap resamples institutions but holds the published Wapman prestige ranks fixed, as if the academic-reputation axis were known without error. If AR ranks carry sampling uncertainty, propagating it should widen the gap CIs and could, in principle, push a ``reliable'' field back inside the noise band, flipping its reliability flag. We inject AR-rank noise by multinomially resampling each field's faculty-placement edges, re-fitting canonical SpringRank on each draw, and using the per-institution percentile-rank SD to perturb ranks \emph{inside} the gap bootstrap. The diagnostic compares the AR-fixed baseline CI width to the AR-propagated full CI width.

\[
\text{median}\ \frac{w_{\text{full}}}{w_{\text{base}}} = 1.04 \ (\text{mean }1.05),\qquad \text{reliability flips} = 0/16 .
\]

\paragraph{Reading.} $w_{\text{base}}$ is the gap-CI width with AR ranks fixed (institution resampling only); $w_{\text{full}}$ adds the AR-rank perturbation from edge-bootstrapped SpringRank; the ratio is taken per field and we report its median over the 16 reliable fields. A ``flip'' is a field whose AR-propagated CI now overlaps the noise band, i.e.\ a field that \emph{loses} its reliable status once AR uncertainty is admitted. We chose to propagate AR uncertainty through the \emph{edges}---resampling placements and re-running SpringRank---rather than through a parametric jitter on the ranks, because edges are the actual data-generating object in the Zenodo release, so the induced rank perturbation respects the network's structure (a top department's rank is more stable than a middling one's). What this \emph{buys} is an honest, structure-respecting widening of the gap CIs; what it \emph{assumes}, and this is the binding caveat, is that the edge-bootstrap SD captures the \emph{relevant} AR uncertainty.

The decode of $1.04$ and $0/16$: propagating AR-rank sampling uncertainty widens the gap CIs by only about 4\% at the median, and flips \emph{no} reliability flags. Substantively, the prestige axis is stable enough that admitting its sampling error barely perturbs the gap inferences---the reliability screen is robust to AR uncertainty. But the crucial interpretive point is that $1.04$ is a \emph{lower bound} on the true cost. The edge bootstrap measures finite-placement sampling noise on the \emph{public, aggregated} edge lists; it does \emph{not} measure the gap between those public edges and the proprietary AARC census that Wapman used---the very gap that caps the from-scratch SpringRank reproduction at $\rho\approx0.77$--$0.8$. So $1.04$ certifies that \emph{edge-sampling} noise on the public network is cheap; it does not certify that the public-vs-census data gap is cheap. We therefore read individual field CIs as somewhat wider than even the AR-propagated bootstrap implies, while noting that the qualitative reliability structure---which fields clear the screen---holds. The per-field table makes the mechanism visible: well-populated, high-prestige-stable fields like nursing ($1.00$) and mathematics ($1.01$) barely move, while thinner or noisier fields like economics ($1.18$) and finance ($1.17$) absorb most of the widening---exactly where AR-rank uncertainty should bite hardest, and still not enough to flip a single flag.

\subsection{The crosswalk audit: is the field-count ceiling a bug or a fact?}

\paragraph{Threat: the analysis ran on only 54 fields because the hand-built Wapman$\to$CIP crosswalk was naive, so the headline is an artefact of an arbitrarily truncated field universe.}
Wapman's taxonomy has 81 fields; the analysis pipeline used 54. A referee can reasonably suspect that the missing 27 are a crosswalk failure---valid undergraduate fields that a better map would have recovered---and that recovering them would move the headline. The audit classifies all 27 missing fields by \emph{why} they are missing, then re-runs the entire decomposition on the recovered, expanded set.

The 27 missing fields partition as:
\[
\underbrace{12}_{\text{recoverable}} \;+\; \underbrace{3}_{\text{subsumed}} \;+\; \underbrace{12}_{\text{genuine gaps}} \;=\; 27 .
\]

\paragraph{Reading.} ``Recoverable'' means a valid undergraduate CIP with Scorecard earnings on $\ge10$ institutions existed but was absent from the hand-built map (e.g.\ Linguistics, Nutrition Sciences, Exercise Science/Kinesiology, Theological Studies)---a genuine crosswalk miss. ``Subsumed'' means the Wapman field is a finer split already contained in an aggregate field, not a gap (Cell Biology$\to$biology, Natural Resources$\to$environmental\_sciences, Information Science$\to$computer\_science). ``Genuine data gap'' means there is \emph{no clean bachelor's-in-field earnings cell}: the field is graduate-only (Epidemiology, Pathology, Pharmacology, Veterinary Medical Sciences, Educational Psychology, $\ldots$) or its bachelor's cell is suppressed at small $n$ (Classics $n=3$, Soil Science $n=6$, Curriculum \& Instruction $n=4$). This three-way partition is the right diagnostic because it separates a \emph{fixable} mapping error from an \emph{irreducible} property of the open data---only the first kind threatens the headline. What it \emph{buys} is an honest accounting of the field-count ceiling; what it \emph{assumes} is that the recovered CIPs match the Wapman PhD field they are joined to (the disclosed taxonomy-slippage risk---e.g.\ Spanish leaning on the broad Romance-languages CIP).

The decode: the field ceiling is \emph{mostly genuine}, not a naive crosswalk failure. Only 12 of 27 are recoverable; the plurality (another 12) are genuine data gaps---disciplines where the academy ranks PhD programs but no bachelor's earnings exist, so no gap is even definable. Recovering the 12 lifts the universe $54\to66$ fields ($48\to57$ non-degenerate, $16\to20$ reliable units). Re-running the decomposition holds the headline on every axis:
\[
\rho_{\text{rank}}(\text{base},\text{expanded}) = +0.98, \quad
\mathrm{ICC}_{\text{raw}}: 0.30 \to 0.32, \quad
\mathrm{ICC}_{\text{corr}}: 0.45 \to 0.57, \quad
\tau^2: 0.0231 \to 0.0207 .
\]

\paragraph{Reading.} The cluster-ranking Spearman $+0.98$ (over shared clusters) says the integrated--decoupled order is essentially identical after expansion---the between-discipline structure is not an artefact of the original 54. The raw ICC barely moves ($0.30\to0.32$): adding 12 fields, distributed across new and existing clusters, leaves the raw between-share almost unchanged. The corrected ICC rises more, $0.45\to0.57$, and this is interpretable: the expanded set adds clusters and fields whose between-cluster signal survives noise-correction better, while $\tau^2$ ticks slightly \emph{down} ($0.0231\to0.0207$) because the expanded, more balanced cluster composition pools more tightly---so the corrected denominator's within-term shrinks proportionally more than the numerator, pushing the corrected ratio up. The key substantive reading is that the headline \emph{survives universe expansion} and the recovery is honestly \emph{capped}: the 12 genuine-data-gap fields are not forced in (recovery would create count, never the missing undergraduate earnings), and the per-institution CIs are unchanged---recovery is a count/cluster-precision gain, not an earnings-precision gain. Like the measurement correction, this check \emph{strengthens} rather than merely survives: the between-discipline share is higher on the expanded universe, and the bulk of the field ceiling is certified as a property of the data, not of our map.

\subsection{Cross-source replication: an independent earnings definition}

\paragraph{Threat: the gap is an artefact of College Scorecard's specific earnings construction (Title-IV-only, median, single cross-section).}
The most serious external threat is that the entire gap rides on idiosyncrasies of one earnings source. The decisive test is to swap in a \emph{completely different} earnings definition---Census PSEO, built from state UI wage records, population-inclusive of UI-employed graduates, in real 2023 dollars---and ask whether both the institution pay ranking and the field-level gap reproduce. Critically, the two sources are \emph{never merged}: PSEO and Scorecard have different populations and PSEO's coalition coverage is non-random, so we compare only the \emph{within-source rank gap} across them.

\[
\rho\bigl(y_{\text{PSEO}},\,y_{\text{Scorecard}}\bigr)=+0.94\ (n=1535),\qquad
\rho\bigl(\text{gap}_{\text{PSEO}},\,\text{gap}_{\text{Scorecard}}\bigr)=+0.68\ [+0.52,+0.80]\ (57\text{ fields}).
\]

\paragraph{Reading.} The first Spearman is computed on the same institution$\times$field cells: do UI wages and Title-IV medians \emph{rank} institutions' pay similarly? The second---the actual test---asks whether the \emph{field-level gap}, recomputed entirely inside PSEO, agrees with the gap recomputed entirely inside Scorecard, across the 57 fields with enough PSEO$\cap$Wapman institutions. We test at the level of the rank gap, not the dollars, because the two definitions are not the same number (UI all-employed vs Title-IV median); comparing gaps---each a within-source rank statistic---is the only join that is immune to the level difference. What it \emph{buys} is the strongest possible robustness claim, replication under an independent population definition; what it \emph{assumes} is that the 108-institution PSEO$\cap$Wapman overlap is informative despite being coalition-public-skewed (elite privates, the brand tail, are under-represented---exactly where Chetty's tail effects live).

The decode: $+0.94$ at the level certifies that the two earnings constructions, despite different populations and dollar definitions, rank institution pay almost identically---they are co-moving measurements of the same underlying earnings, not the same number. $+0.68$ on the gap is the headline-bearing figure: an \emph{independent} earnings source reproduces not just the levels but the prestige$\leftrightarrow$pay \emph{misalignment} structure, with a CI $[+0.52,+0.80]$ comfortably clear of zero. The gap is below the level agreement ($0.68<0.94$) for an instructive reason---the gap is a second-order rank statistic (a function of a within-field correlation), so it inherits and amplifies the residual disagreement between the two sources, and the 57-field overlap is thinner and more coalition-skewed than the 1535-cell level comparison. So $+0.68$ is exactly the attenuated-but-decisive value one expects when a derived rank statistic replicates across genuinely independent inputs.

\paragraph{Geography as an explicit, measured control.}
PSEO's destination flows let us turn the geography confound from an assumption into a measured control.
\[
\text{out-of-state at 5yr}=34\%,\qquad R^2_{\text{origin-state}}\approx0.13,\qquad \overline{\text{gap}}_{\text{PSEO}}:\ 0.68 \to 0.67 .
\]

\paragraph{Reading.} The $34\%$ out-of-state share (i.e.\ only 66\% employed in the institution's state at 5 years) is both a substantive mobility fact and a \emph{bound} on the PSEO undercount: a single state's UI files miss the third of graduates who leave, so PSEO's within-state earnings are an undercount whose magnitude this number caps. $R^2\approx0.13$ is the share of within-field cross-institution earnings explained by origin-state fixed effects on this overlap. The last quantity is the mean PSEO gap before and after residualising earnings on destination-state fixed effects. We use destination-state FE rather than full wage modelling because state is the geographic unit PSEO releases and the question is narrow---does \emph{where graduates land} drive the gap? What it \emph{buys} is an explicit geography control; what it \emph{assumes} is that destination-state FE captures the relevant geography (selection/value-added remains unaddressed---a high-paying destination may reflect who moves there).

The decode of the apparent tension: geography drives earnings \emph{levels} ($R^2\approx0.13$ here, and a larger $\approx0.46$ on Scorecard institution-mean residuals in the level analysis), yet \emph{netting geography barely moves the gap} ($0.68\to0.67$, a one-point shift over 60 fields). This is the load-bearing interpretation. Geography is a real component of where institutions' graduates earn, but the gap is a \emph{within-field rank} object: shifting every institution's earnings by a common origin- or destination-state effect rescales levels without reordering the within-field prestige$\leftrightarrow$pay ranking, so the rank gap is nearly invariant to it. The $0.13$ being modest \emph{on this overlap}---smaller than the $0.46$ elsewhere---is itself decoded in the source: the 108 PSEO$\cap$Wapman institutions are mostly coalition publics, a narrower and more homogeneous sample, and the destination measure is coarse (institution-level, all-CIP), both of which compress geography's measured share. The honest verdict the number certifies: destination geography moves earnings \emph{levels/residuals} but not the field-level prestige$\leftrightarrow$pay \emph{rank gap}.

\subsection{Multi-horizon stability: is the gap a quirk of the 4-year window?}

\paragraph{Threat: the 4-year-post-completion horizon is arbitrary, and the gap (and its discipline structure) would look different at another horizon.}
We re-estimate the gap and the CIP-2 decomposition at every consistent Scorecard horizon---1, 4, and 5 years post-completion---holding AR (Wapman SpringRank) fixed and changing only ER. The diagnostic is whether the ICC level and, more importantly, the cluster and field rankings are stable across horizons.
\[
\mathrm{ICC}:\ 0.50\ (\text{1yr}) \to 0.30\text{--}0.33\ (\text{4--5yr}),\quad
\rho_{\text{cluster}}=+0.81\text{ to }+0.85,\quad
\rho_{\text{field-gap}}=+0.68\text{ to }+0.74 .
\]

\paragraph{Reading.} The ICC at each horizon is the same between-cluster share statistic recomputed on that horizon's earnings; $\rho_{\text{cluster}}$ and $\rho_{\text{field-gap}}$ are the pairwise Spearman rank correlations of cluster gaps and field gaps across the three horizon pairs (1-vs-4, 1-vs-5, 4-vs-5). We test \emph{ranking} stability separately from the \emph{ICC level} because the two answer different questions---whether the integrated--decoupled \emph{ordering} is horizon-robust (it is) versus whether the magnitude of the between-share is horizon-constant (it is not). What this \emph{buys} is scope clarity; what it \emph{assumes} is that 1--5 years is the relevant window---we explicitly do not extrapolate to lifetime earnings, where late-blooming fields (medicine-adjacent, law-adjacent) could re-rank.

The decode of the ICC drift: the ICC is \emph{not} flat---it is $0.50$ at 1 year and settles to $0.30$--$0.33$ at 4--5 years. This is honest and interpretable, not a contradiction of the headline. One year out, within-field earnings are barely differentiated (everyone starts near the same salary), so gaps are large and \emph{more} of their variance sits between disciplines; the grand gap is correspondingly higher at 1yr ($0.77$) than at 4--5yr ($\sim0.56$). As employers learn and within-field pay disperses over the first few years, the between-discipline share \emph{declines} toward the $0.30$--$0.33$ that the paper reports as its headline (the 4--5yr figure). So the decline is an \emph{early-career} learning phenomenon, and it moves toward \emph{more} between-discipline structure at shorter horizons---the headline direction strengthens, not weakens, earlier. The ranking stability ($\rho_{\text{cluster}}=+0.81$ to $+0.85$, $\rho_{\text{field-gap}}=+0.68$ to $+0.74$) certifies that the integrated end (Computing, Math \& Stats, Social Sci, Engineering) and the decoupled end (Health, natural-science and arts clusters) keep their places across horizons. The essential caveat, which makes this a \emph{necessary-but-weak} test, is that the horizons are \emph{autocorrelated}: the 1/4/5-year figures are the same IRS-linked completers re-measured, so cross-horizon stability is partly mechanical. Instability \emph{would} have falsified the headline; observed stability is consistent with it but does not independently prove it---which is precisely why this check is supported by the genuinely independent PSEO cross-source test above.

\subsection{Field-tag robustness: research field is not degree field}

\paragraph{Threat: the gap's residual to the data ceiling is an artefact of noisy free-text field assignment, fixable by better tagging.}
The prestige axis is built by assigning each faculty placement to a field via a degree/department text string, which is noisy. If that noise is what holds per-field reproduction of Wapman below the $\sim0.79$ public-Wapman ceiling, then a principled field tagger should close the gap. We replace the free-text assignment with OpenAlex publication-topic \emph{research}-field tags and re-run the per-field SpringRank-vs-Wapman validation head-to-head.
\[
\overline{\rho}_{\text{OpenAlex}}=0.515 \;\;\text{vs}\;\; \overline{\rho}_{\text{degree-text}}=0.551;\qquad
\text{de-leaked: } 0.515 \;\;\text{vs}\;\; 0.517 .
\]

\paragraph{Reading.} Each $\rho$ is the mean across cleanly 1:1-mappable fields of the Spearman between a SpringRank fit on that tagging scheme's subnetwork and Wapman's published field ranks; the benchmark ceiling is $\overline{\rho}\approx0.791$. The ``de-leaked'' pair re-runs the degree-text baseline after dropping rows whose text \emph{also} names a competing field (e.g.\ ``computer'' catching electrical-\&-computer engineering, ``psycholog'' catching educational psychology), which had inflated the raw degree-text score. We compare at \emph{matched edge count} and after de-leaking because the raw $-0.036$ OpenAlex deficit is otherwise misleading---it conflates tagging quality with coverage. What this \emph{buys} is a clean separation of \emph{tag-noise} from \emph{coverage} as the binding constraint; what it \emph{assumes} is that the cleanly-mappable fields are representative.

The decode: OpenAlex research-field tags do \emph{not} beat degree-text ($0.515<0.551$ raw), and---decisively---once the degree-text baseline is de-leaked the two are a \emph{dead heat} ($0.515$ vs $0.517$). The raw gap was almost entirely (i) keyword leakage inflating the degree-text baseline and (ii) one pathological field, Chemistry (excluding it, the mean $\Delta$ is $-0.014$, a wash). The substantive reason both schemes plateau at $\sim0.515$ rather than reaching $0.79$ is that \emph{research field is not degree field}: only 35--54\% of a degree cohort (Chemistry just 11\%) lands in the matching OpenAlex research field---chemists publish into biochemistry and materials, psychologists into social sciences and medicine---so research-field tagging scatters each degree cohort across adjacent fields and \emph{thins} every per-field subnetwork. The reproduction ceiling is therefore \emph{coverage-bound, not tagging-noise-bound}: neither scheme reaches the benchmark even at matched $n$, so the residual to the $\sim0.79$ ceiling reflects sparse public ORCID edges against Wapman's proprietary AARC census, and \emph{better labels cannot close it---only more edges can}. This has a constructive corollary for the headline: because the two taxonomies cross-cut at the fine-field level but converge at the coarse discipline-cluster level, the cluster grain is exactly where the AR--ER gap structure is robust (ICC$\approx0.30$)---the field-tag check reinforces the choice to read the gap at the between-discipline-cluster level rather than chase fine-field precision the open data cannot support.\section{Temporal revaluation and applied extensions (decoded)}\label{sec:temporal-applied}

The static spine of the paper treats the gap as a single cross-section: prestige and pay are two orderings frozen at one moment. This section explains the three analyses that push past that frozen frame --- one that adds a \emph{time} axis to the earnings side (revaluation), one that asks whether the prestige side moves to track it (the lead--lag pilot), and one that turns the field-specific structure into an \emph{applied artifact} (the CS-school ranking). Their epistemic statuses are deliberately unequal, and the whole point of decoding them is to keep those statuses visible: the revaluation is a clean novel descriptive fact, the lead--lag is suggestive but unidentified, and the ranking is an engineering demonstrator that illustrates the paper's critique rather than testing a scientific claim. We take them in that order.

\subsection{Temporal revaluation: the labour market actively reprices fields}\label{sec:reval}

\paragraph{Design rationale.} Everything upstream measures $\mathrm{gap}_f$ on a single Scorecard cross-section, so a natural objection is that the prestige--pay disagreement is a fixed property of fields --- a snapshot artifact that says nothing about whether the two valuation systems are even pointed at a moving target. The revaluation analysis answers the precondition for that objection: \emph{is the labour-market side moving at all?} It uses Census PSEO field-level bachelor's earnings (5-year-after-completion, cohort-weighted) across graduation cohorts 2001--2016, and crucially these are \emph{real} 2023 dollars --- already inflation-adjusted in the LEHD schema --- so any trend is a genuine repricing of the field, not a price-level drift. For each field with at least four cohort points (60 fields qualify) the analysis fits a per-field linear trend in real dollars per year. The estimand is a slope, not a level, so it is read as ``how many constant-2023 dollars per year is this field's median graduate gaining or losing.''

\paragraph{The rising fields, decoded.} Six fields anchor the upward end. Statistics rises $+\$2{,}311$/yr (from $\$84{,}447$ in 2007 to $\$102{,}443$ in 2016), Computer Science $+\$1{,}616$/yr ($\$85{,}674\to\$107{,}600$, 2001--2016), Computer Engineering $+\$1{,}421$/yr ($\$104{,}037\to\$122{,}368$), Biomedical Engineering $+\$969$/yr ($\$79{,}406\to\$92{,}273$), Finance $+\$868$/yr ($\$80{,}486\to\$91{,}189$), and Economics $+\$817$/yr ($\$76{,}829\to\$86{,}774$). These are not arbitrary: every one of them is a field that the static decomposition placed at or near the \emph{integrated} (low-gap) end of the typology --- Computing $0.348$, Math \& Stats $0.352$, Economics-adjacent social science $0.421$, Engineering $0.432$ in \S4. The substantive reading is that the fields the labour market prices most coherently against prestige are also the fields it is actively bidding \emph{up} in real terms, which is exactly what one would expect if these are the fields where market demand and observable skill are both rising together. The magnitudes are economically large: $+\$2{,}311$/yr compounded over the nine-cohort Statistics window is the $\sim\!\$18{,}000$ real gain that the first/last figures show, i.e.\ the median statistics bachelor's graduate of 2016 out-earned their 2007 counterpart by roughly a fifth, in constant dollars.

\paragraph{The falling fields, decoded.} The two sharp declines are Pharmacy $-\$1{,}774$/yr ($\$150{,}946\to\$119{,}777$) and Communication Disorders $-\$1{,}405$/yr ($\$95{,}936\to\$74{,}500$). These take their values for a specific and informative reason: both are \emph{licensed, supply-regulated} health fields, exactly the kind of field that the licensing channel (\S5, $b_{\text{licensure}}=+0.66$) marks as compressed. Pharmacy in particular underwent a well-documented expansion in the number of accredited PharmD programs over precisely this window, and a real decline of $\sim\!\$31{,}000$ over fifteen cohorts is the labour-market signature of that supply shock: when a licensed field floods its own entry, the real median falls even though the credential barrier stays high. The contrast with the rising computing/quantitative fields is the substantive payload --- the market is not uniformly inflating all degrees; it is \emph{reallocating}, lifting some fields and discounting others in real terms. The bulk of fields drift mildly downward or flat (Teacher Ed $-\$68$, Chemical Engineering $-\$81$, English $-\$81$, Music $-\$83$, Mechanical Engineering $-\$136$), confirming that the large movers are genuine outliers, not noise.

\paragraph{The headcount and its meaning.} The summary statistic is \textbf{46/60 fields rising} (14 falling). This number is the load-bearing one: it decodes to ``the modal field is being repriced \emph{up} in real dollars, but a substantial minority is being repriced down,'' which establishes both halves of the claim --- the labour market is dynamic (not a static cross-section), and it is selective (not a rising tide). The honest caveat that travels with every number here is \emph{coalition-composition drift}: PSEO's field series are built from whichever institutions are in the disclosure-released coalition in each cohort, and that coalition's membership shifts over 2001--2019. If a field's covered institutions change composition over time, the \emph{level} of its series can drift for reasons unrelated to true repricing. This is why the result is framed as a within-field trend in real dollars and not as a claim about absolute field rankings: the caveat bites on the levels, far less on the direction of within-field movement, and the 46/60 count is robust because it is a sign count, not a level comparison across fields.

\subsection{The lead--lag pilot: suggestive, deliberately not identified}\label{sec:leadlag}

\paragraph{Why a pilot and not a result.} The exciting question the revaluation opens is the \emph{lead--lag}: if the market reprices a field, does academic prestige adjust to follow, or does prestige sit as a slow anchor while pay drifts around it? Answering that requires a prestige \emph{panel} --- multiple time-stamped reconstructions of the faculty-hiring hierarchy --- and the binding constraint is the prestige side, not the earnings side. Academic prestige is slow-moving, and only two ORCID-rebuilt hiring-network periods (2011--2015 vs.\ 2016--2020) are realistically buildable from the cached edges; each five-year window per field carries only $\sim\!300$--$1500$ edges, so each SpringRank reconstruction is itself coarse. The pilot therefore does not attempt a cross-lagged estimate. It computes one diagnostic --- the within-field period-to-period rank \emph{stability} of prestige --- and uses it to \emph{size} the real study rather than to deliver a finding.

\paragraph{The stability number, decoded.} Across 10 well-covered fields the within-field prestige stability is \textbf{median $\rho\approx0.50$}, ranging $+0.34$ (Physics) to $+0.58$ (Economics), with Computer Science $+0.57$, Mathematics $+0.53$, Psychology $+0.53$, Mechanical Engineering $+0.55$, Chemistry $+0.46$, Electrical Engineering $+0.47$, Biology and Civil Engineering both $+0.45$. The single most important interpretive move is that $\rho\approx0.50$ is a \emph{noise-depressed lower bound on true stability}, not an estimate of it. Two facts force this reading. First, each period's prestige is built on only $300$--$1500$ ORCID edges, so the two reconstructions being correlated are each individually noisy, and measurement error in both axes attenuates a correlation toward zero. Second, even the full public Wapman SpringRank reproduces \emph{itself} only at $\rho\approx0.77$ (the data-ceiling reproduction from \S3); a network that cannot be perfectly re-derived from its own richer data certainly cannot show period-to-period stability above that ceiling under thinner edges. So the true period-to-period stability of academic prestige is higher than $0.50$ --- $0.50$ is what survives the noise. Read against Part~1, where real earnings move materially over the same window, the juxtaposition is \emph{suggestive} that prestige is the slower-moving anchor and the market is where revaluation happens, which is directionally consistent with the structural model's prediction P5 (academic value as a slow prior, earnings drifting).

\paragraph{Why the direction is genuinely not identified.} The decode that must not be softened: \textbf{the cross-lagged direction is NOT identified with two noisy periods}. A real lead--lag test regresses $\Delta\text{prestige}_t$ on $\Delta\text{earnings}_{t-1}$ and, symmetrically, $\Delta\text{earnings}_t$ on $\Delta\text{prestige}_{t-1}$, and reads the lead--lag off the sign/magnitude contrast --- but a contrast of two \emph{changes} requires at least three waves, and the pilot has two. With two waves there is exactly one prestige change and one earnings change per field; ``which one leads'' is mathematically undefined, not merely imprecise. The stability \emph{asymmetry} (prestige moderately stable while earnings move) is the most one can honestly report, and it is a directional hint, not a causal estimate that academia provably lags. The specification for the real study is therefore stated and not run: rebuild ORCID SpringRank in $\ge4$ rolling 3-year windows and run the cross-lagged panel. The second, deeper data block is recorded honestly: the \emph{institution-level} earnings panel that a within-institution cross-lagged design would need is unavailable because PSEO single-cohort institution$\times$field earnings cells are almost entirely disclosure-suppressed --- the same suppression wall (recall doctoral cells at $0/18{,}442$ released) that bounds the rest of the paper. The pilot thus demonstrates feasibility and prices the real study; it does not deliver the lead--lag, and we are explicit that it cannot.

\subsection{The CS-school ranking: an applied demonstrator of the field-specific critique}\label{sec:csrank}

\paragraph{Why build it at all.} The paper's practical implication is a critique of field-agnostic prestige rankings: if the prestige$\leftrightarrow$pay map is field-specific (the integrated--decoupled typology), then a single overall brand ranking mis-prices programs. The CS-school ranking makes that critique concrete by building the thing a field-agnostic ranking \emph{cannot} produce --- a ranking internal to one field that fuses that field's academic prestige with that field's labour-market pay. Its status is an applied demonstrator, not the scientific contribution; we decode its numbers fully precisely so that the descriptive limits are unambiguous.

\paragraph{The blend, presented.}
\[
\text{combined}_i \;=\; w\,z_{P,i} \;+\; (1-w)\,z_{S,i}, \qquad w = 0.6 .
\]

\paragraph{Reading.} Here $z_{P,i}$ is school $i$'s \emph{standardized CS prestige} --- its continuous SpringRank score in the CS faculty-hiring subgraph, recomputed from the public CS edge subgraph so that score \emph{differences} (not just ranks) are meaningful, then $z$-standardized to mean $0$ and SD $1$ across the 70-school CS universe. $z_{S,i}$ is school $i$'s \emph{standardized CS salary} --- its College Scorecard CS bachelor's median earnings (4-year-after-completion), likewise $z$-standardized over the same universe. The weight $w=0.6$ is the convex mixing parameter: $60\%$ of the combined score is prestige, $40\%$ is salary, so the ranking is \emph{prestige-anchored} with salary as a fine-tuning correction. Three design choices are encoded in this one line. (i) Why $z$-standardize before blending: prestige is a dimensionless SpringRank score of order one while salary is in dollars of order $10^5$; adding the raw quantities would let dollars swamp prestige purely through units, so each axis is first put on a common, unit-free scale (SD units) before they are mixed --- the standardization is what makes $w$ a genuine relative weight rather than an artifact of measurement units. (ii) Why a \emph{convex} blend ($w$ and $1-w$ summing to one) rather than a free two-coefficient regression: the goal is a transparent, interpretable ranking knob, not a fitted model --- a single $w$ that a reader can dial lets the sensitivity of every school be read directly off how its rank moves with $w$, which is the basis for the robustness flagging below. (iii) What the blend \emph{buys} and \emph{assumes}: it buys a field-specific ordering that a field-agnostic brand ranking structurally cannot deliver (because the latter has no field axis at all); it \emph{assumes} that prestige and pay are commensurable once standardized and that the prestige anchor ($w>0.5$) is the more trustworthy axis, with salary only fine-tuning --- an assumption the caveats below make explicit by treating salary-driven movements as the \emph{less} robust ones.

\paragraph{The top of the table, decoded.} At $w=0.6$ the order is Carnegie Mellon, Stanford, Berkeley, MIT, UCLA. CMU lands first with a CS median of \textbf{\$246{,}966} and a \emph{salary premium} of \textbf{$+\$48{,}910$}. The premium must be decoded carefully because it is not the raw salary: it is the residual from regressing CS earnings on CS prestige score (slope $\$68{,}449$ per SpringRank unit), so $+\$48{,}910$ means CMU's CS graduates earn about \$49k \emph{above what its academic standing alone predicts}. That positive residual is exactly why CMU tops a blend that an academic-only ranking would not place first: its salary axis pulls it up past schools of comparable or higher pure prestige. The $w$-sensitivity columns decode the anchoring directly --- CMU is rank 1 at $w=0.5$ but slips to rank 3 at $w=0.7$, because as prestige weight rises its salary advantage is down-weighted; symmetrically Berkeley (a $-\$6{,}379$ premium, i.e.\ paid slightly \emph{below} its standing) climbs from rank 4 at $w=0.5$ to rank 1 at $w=0.7$. This is the demonstrator working as intended: the salary axis is doing visible, interpretable work at the top.

\paragraph{The two cross-validations, decoded.} The ranking's two axes are each independently cross-checked. The \emph{earnings} axis: Scorecard CS earnings against Census PSEO CS earnings on the $n=29$ overlap gives Spearman \textbf{$+0.84$}. The \emph{prestige} axis: Wapman CS SpringRank against the ORCID-rebuilt CS SpringRank on the $n=61$ overlap gives Spearman \textbf{$+0.69$}. Both numbers are below $1.0$ and both are below their corresponding cross-validations elsewhere in the paper (e.g.\ the institution-level Scorecard$\leftrightarrow$PSEO pay correlation of $+0.94$), and the decode is that this is a \emph{data ceiling, not a defect}. The $+0.84$ earnings agreement is lower than the $+0.94$ institution-level figure because it is computed on a small ($n=29$), CS-only overlap where two genuinely different populations (Title-IV recipients vs.\ UI-employed graduates) are being compared field-by-field; a Spearman in the mid-$0.8$s across two independent earnings \emph{definitions} on 29 schools is strong agreement, not weak. The $+0.69$ prestige agreement is bounded above by the very reconstruction ceilings established earlier --- the public Wapman network reproduces itself only at $\rho\approx0.77$ and the ORCID rebuild matches Wapman academia ranks at $\rho=0.74$ overall --- so a per-field CS $\rho=+0.69$ on thin CS-only edges sits right at the achievable ceiling and \emph{certifies} that both axes replicate on data the primary build never touched. The substantive payload: neither axis of the ranking is an artifact of a single source.

\paragraph{The robustness flags, decoded.} \textbf{15 schools} carry a low-robustness flag, defined operationally as a rank that moves more than 5 positions between $w=0.5$ and $w=0.7$ --- i.e.\ schools whose placement swings materially with how much weight you give salary. These are not random: they are exactly the \emph{salary-driven} entries (USC, Ohio State, Washington University in St.\ Louis, George Washington, UMass Amherst, Chicago, Dartmouth, Cincinnati, Yale, Michigan, Maryland, Tufts, Penn State, Kent State, UC Davis). The decode is the cleanest possible vindication of the field-specific critique: these are precisely the programs whose rank depends on the field-specific pay axis, so they are precisely the entries a field-agnostic brand ranking --- which has \emph{only} the prestige axis --- would place most wrongly. Dartmouth is the textbook case: a $+\$79{,}328$ salary premium drags it from rank 22 at $w=0.5$ down to rank 39 at $w=0.7$ as the salary boost is stripped away. The flag is therefore not an admission of weakness in the demonstrator but its core illustration --- it points at the very schools where ``field-agnostic vs.\ field-specific'' makes the largest difference.

\paragraph{The descriptive boundary, decoded.} The last number is the geography check: destination-state geography explains \textbf{$R^2\approx0.13$} of the within-field earnings variation here (the same origin-state figure that bounds the gap robustness in \S7). This is the load-bearing caveat for reading the salary axis. It decodes to: the salary term reflects \emph{where graduates land}, not a school's value-added ceiling. Roughly $34\%$ of bachelor's graduates work out-of-state, so a strong school whose graduates concentrate in a lower-wage region is ``dinged'' on the median even though a mobile graduate from it might do fine --- this is why several Midwest and regional flagships sit below coastal schools of similar prestige. The median is a placement statistic, not a productivity statistic. Combined with the explicit non-causal stance --- a high salary premium is \emph{not} ``this school makes you richer,'' because \citet{chetty2023diversifying} locate the elite-brand causal payoff in the top tail that median data cannot see, and selection is unaddressed --- the ranking is honestly bounded as \emph{descriptive}. It is built on the Title-IV Scorecard population, but the $+0.84$/$+0.94$ cross-validations show the \emph{ordering} survives the population definition. In sum: the ranking demonstrates the paper's practical critique --- a field-specific blend reveals exactly the salary-driven, low-robustness entries a field-agnostic ranking would misprice --- without overstepping into a causal value-added claim it cannot support.\section{The measurement asymmetry: earnings as a proxy for employer reputation}\label{sec:asymmetry}

\subsection{The thesis: the two sides of the gap are not measured with comparable richness}

The central statistic of this paper, $\mathrm{gap}_f = 1 - \mathrm{Spearman}_i(\mathrm{prestige}_{i,f},\,\mathrm{earnings}_{i,f})$, is a difference between two orderings. Every interpretive weight we have placed on it --- ``integrated versus decoupled disciplines,'' ``valuation divergence between the academy and the market,'' ``mis-priced programs'' --- presupposes that the two orderings are commensurable readings of the same kind of object: \emph{where the academy ranks a program} versus \emph{where the labour market ranks it}. This section argues that they are not, and that the asymmetry is not a footnote but the deepest interpretive constraint on the whole construction.

The academic-reputation side (AR) is a rich relational object. It is a continuous SpringRank score \citep{debacco2018physical} extracted from the entire US faculty-hiring network \citep{wapman2022quantifying}: a hierarchy inferred from thousands of placement edges, where each edge --- department $j$ hires department $k$'s PhD --- is a revealed pairwise judgement of standing, and the score is the latent rank that best rationalises the whole tournament of those judgements simultaneously. AR is therefore a \emph{network} measurement: it aggregates a high-dimensional web of relational, peer-revealed preferences into a score, and it does so over a population (faculty hires) that is itself a curated, expert sample of the academy's own valuation.

The employer-reputation side (ER), as actually operationalised, is a single scalar: median field-of-study earnings at one horizon \citep{scorecard}. Against a network of thousands of hiring decisions we set one number per institution-field cell --- the central tendency of one outcome (wages) of one slice of graduates (Title-IV) at one moment in the career (roughly four years out). The two sides of the wedge are not measured with comparable richness. AR is a relational census of a community's revealed esteem; ER is a point estimate of one moment of one outcome. The gap is a comparison across a measurement-richness gradient, and any reading of its residual must be disciplined by that fact.

\subsection{In fairness to earnings: why the median wage is the most defensible scalar}

Before pressing the criticism, the scalar deserves its defence, because the choice of median earnings was not lazy --- it was the single most defensible reduction available. The labour market does not publish a ``reputation rank'' the way the academy implicitly does through hiring. What it publishes, through the wage, is a \emph{revealed price}. In the signalling and employer-learning frame of \S2 \citep{spence1973job, altonji2001employer}, the wage is precisely the market's posterior valuation of a worker's productivity after it has folded in every signal it can observe --- credential, brand, demonstrated skill, occupational sorting. The median wage is thus not an arbitrary coordinate of ER; it is the one coordinate into which the market itself compresses all the others. Beyond being a revealed price and (below) cross-validated, it has four further properties no competing scalar matches at this scale: it is \emph{comparable} across every field and institution on one ratio scale; it is \emph{administratively grounded}, anchored to IRS/UI tax records rather than self-report; it is \emph{available} at the institution-by-field grain the gap requires; and it is the outcome that students and policy-makers most directly care about and act on.

Crucially, it is also \emph{cross-validated}. This is where the asymmetry argument must concede real ground to the scalar, and the concession is quantitative:
\[
\mathrm{Spearman}\big(y^{\text{PSEO}}_{i,f},\, y^{\text{Scorecard}}_{i,f}\big) = +0.94 \quad (n = 1535),
\]
the level agreement between two genuinely independent earnings definitions on the same institution-field cells \citep{pseo, scorecard}.

\paragraph{Reading.} Here $y^{\text{PSEO}}_{i,f}$ is the Census PSEO UI-wage-based earnings figure for cell $(i,f)$ --- built from state unemployment-insurance wage records, in real 2023 dollars, covering UI-employed graduates --- and $y^{\text{Scorecard}}_{i,f}$ is the College Scorecard Title-IV median for the same cell. The Spearman (rather than Pearson) coefficient is the right summary because the gap is itself a rank statistic: what must be stable for our construction to be meaningful is the \emph{ordering} institutions induce on pay, not the dollar levels, which legitimately differ between a Title-IV median and an all-UI-employed mean. The value $+0.94$ is high but pointedly \emph{not} $1.00$, and the residual $0.06$ is substantively the right size: it is not noise but the genuine population difference between ``federally-aided completers'' and ``UI-employed graduates'' (the latter misses the federally-employed, the self-employed, and out-of-state movers with no coalition state). What $+0.94$ buys is the assurance that the median-wage coordinate is not an artifact of one agency's definitional quirks: two bureaucracies, two populations, two dollar conventions, one ordering. What it assumes is that rank-agreement on the \emph{observed} coordinate licenses treating that coordinate as a faithful index of ER --- and it is exactly that assumption the rest of this section interrogates. The $+0.94$ certifies the wage as a robust measurement of \emph{the wage}. It says nothing about whether the wage is a robust measurement of \emph{employer reputation}.

\subsection{The cost: employer reputation is multi-dimensional, and the median collapses it}

``Employer reputation,'' or more carefully a field's labour-market standing, is not a scalar latent. It is a vector. Write it schematically as
\[
\mathrm{ER}_{i,f} = \Big(\underbrace{\text{employer identity/selectivity}}_{d=1},\ \underbrace{\text{placement/underemployment rate}}_{d=2},\ \underbrace{\text{occupational attainment/status}}_{d=3},\ \underbrace{\text{trajectory/growth}}_{d=4},\ \underbrace{\text{non-wage rewards}}_{d=5},\ \underbrace{\text{median wage}}_{d=6}\Big),
\]
of which the analysis observes only the sixth coordinate.

\paragraph{Reading.} Each $d$ is a distinct, non-redundant axis along which a labour market confers or withholds standing. $d=1$ is \emph{which} employers hire the graduates --- the prestige and selectivity of the hiring firms, the labour-market analogue of the faculty-hiring edges that build AR. $d=2$ is the rate of successful placement and its converse, underemployment: whether graduates work at all and at the level their degree implies. $d=3$ is occupational attainment and status --- the surgeon-versus-aide distinction, captured by an occupational-prestige scale, which two fields with identical median pay can rank oppositely. $d=4$ is the \emph{slope} of the earnings path, not its level: a field can have a modest median four years out yet a steep trajectory (or the reverse). $d=5$ is the bundle of non-wage rewards (stability, autonomy, benefits, mission) that compensating-differentials theory says trades off against wages. $d=6$ is the median wage, the coordinate we measure. The functional form here is deliberately a \emph{vector}, not a weighted index, because writing it as a vector makes the loss explicit: any scalar ER is a projection of this vector onto one direction, and median earnings projects onto the basis vector $e_6$ alone. What the vector representation buys is a precise statement of the omission --- we are not measuring a noisy ER, we are measuring one of its six coordinates and silently treating it as the whole. What it assumes (and what is false) is that the six coordinates are collinear enough that $d=6$ proxies for $(d{=}1,\dots,d{=}5)$. They are not: two programs with identical median earnings can have entirely different employer-selectivity, placement, occupational-status, trajectory and non-wage profiles. Median earnings collapses a six-dimensional standing into one central-tendency wage, and the collapse is many-to-one.

This is why ``employer reputation,'' used unqualified, over-claims. The wage is one shadow of ER, the most defensible shadow, but a shadow.

\subsection{The sharp methodological consequence: the residual is confounded with measurement thinness}

The asymmetry has a precise and damaging consequence for the paper's most ambitious reading of the gap. \S6 entertains the interpretation that the licensing-netted residual of the gap estimates \emph{valuation divergence} $(1-\rho_f)$ --- a genuine disagreement between how the academy ($\theta_A$) and the market ($\theta_M$) value the same programs, where $\rho_f = \mathrm{corr}(\theta_A,\theta_M)$. The measurement asymmetry shows this reading is confounded at the source. Consider a field where academic prestige strongly predicts \emph{which employers hire graduates} ($d=1$), or \emph{occupational status} ($d=3$), or \emph{placement} ($d=2$), but \emph{not} the median wage ($d=6$) --- a plausible profile for, say, a field whose elite programs route graduates into prestigious but not top-paying occupations. Our gap, computed only on $d=6$, would read \emph{large} for that field. But the substantive meaning of that large gap is ``wages do not capture this field's reward,'' i.e.\ a thin-measurement artifact --- not ``the academy and the market disagree about quality,'' i.e.\ true divergence $1-\rho_f$. The two are observationally identical in our data and mean opposite things.

So the residual-as-divergence reading is \emph{doubly} uncertain. It is uncertain first for the reason \S6 already concedes: when tested against independent divergence proxies, the make-or-break validation failed. The two proxies gave opposite signs,
\[
\mathrm{Spearman}(\text{residual},\,\text{O*NET}) = +0.41\,[+0.08,+0.68], \qquad \mathrm{Spearman}(\text{residual},\,\text{SDR}) = -0.61\,[-0.88,-0.11],
\]
and were themselves negatively correlated ($-0.64$).

\paragraph{Reading.} The residual is the gap net of the licensing channel, $\mathrm{gap}_f - 0.650\,\lambda_f$ (\S5). The O*NET proxy scores a field by the applied-minus-abstract skill intensity of the occupations its graduates enter; the SDR proxy scores it by the normalised academic-minus-industry salary wedge for doctorate-holders, $|w_{\text{ac}}-w_{\text{ind}}|/\bar w$. The model predicts the residual should correlate \emph{positively} with any honest divergence proxy. O*NET obliges weakly ($+0.41$, CI clear of zero but attenuating to $+0.34$ once field-size is controlled, betraying a size leak); SDR \emph{robustly contradicts} ($-0.61$, CI entirely negative, made \emph{more} negative by licensing-netting, so the contradiction is not a licensure contamination artifact). Two proxies that were genuinely two noisy reads of one divergence construct would correlate positively with \emph{each other}; their correlation of $-0.64$ proves they measure \emph{different things} and disagree. The decisive decoding: this opposite-signs-plus-negative-mutual-correlation pattern is exactly the signature one expects when there is \emph{no single well-defined divergence object} for the residual to be estimating --- which is precisely what the multi-dimensionality of ER predicts. If the residual were partly $d{=}6$-specific measurement thinness rather than pure $1-\rho_f$, then proxies built on different ER dimensions (O*NET leans on occupation/skill content, $d{=}3$; SDR leans on a price wedge, $d{=}6$-adjacent) \emph{should} disagree, because they are projecting the gap onto different ER coordinates. The $+0.41$-versus-$-0.61$ failure is thus not just a failed validation; read through this section it is positive evidence that the residual is not a clean divergence signal. This is why the paper downgrades the residual-as-divergence claim from ``estimated'' to ``unconfirmed.'' The measurement asymmetry supplies the \emph{reason} the validation could not succeed: the residual absorbs ER mis-measurement on top of (or instead of) divergence, so no single external divergence proxy could ever cleanly confirm it. The B3 tie-breaker --- $\mathrm{corr}(\lambda_f,\text{O*NET})=-0.03$ versus $\mathrm{corr}(\lambda_f,\text{SDR})=+0.39$ --- cannot rescue O*NET either: it shows SDR's residual is partly licensure-contaminated but does nothing to establish that O*NET's $+0.41$ is the \emph{right} sign rather than the lucky one.

\subsection{The naming-precision correction}

The honest consequence is a correction of nomenclature that should propagate through the paper. Strictly, the object we measure is the gap between academic prestige and \emph{median earnings} --- the academic-prestige-versus-pay gap. ``Employer reputation'' is the broader latent construct (the six-dimensional ER above) that median earnings \emph{proxies for}. The two are not the same, and the $+0.94$ cross-validation establishes only that the proxy is a stable measurement of pay, not that pay is a stable measurement of ER. Accordingly, wherever the paper labels the second axis ``ER,'' it should be read as ``PAY, as a (defensible but partial) proxy for the broader employer-reputation construct.'' The integrated-versus-decoupled typology of \S4 is, with full precision, an integrated-versus-decoupled typology of \emph{prestige against pay}; whether it is also a typology of prestige against \emph{labour-market standing in general} is an open question this design cannot answer, because it never observed the other five coordinates.

\subsection{The principled fix: a multi-dimensional ER and a per-dimension gap}

The resolution the static paper does not yet have is to stop projecting ER onto a single coordinate. Two routes follow naturally, and both point at the same forward thrust --- bringing in resume-and-employment-trace data (Revelio-style) that observes the missing coordinates. First, one can recompute the gap \emph{per ER dimension}:
\[
\mathrm{gap}^{(d)}_f = 1 - \mathrm{Spearman}_i\!\left(\mathrm{prestige}_{i,f},\ \mathrm{ER}^{(d)}_{i,f}\right), \qquad d \in \{1,\dots,6\}.
\]

\paragraph{Reading.} $\mathrm{ER}^{(d)}_{i,f}$ is the $d$-th coordinate of the ER vector for cell $(i,f)$: $d{=}1$ an employer-selectivity index built from the identities of hiring firms in resume data; $d{=}2$ a placement/underemployment rate; $d{=}3$ an occupational-prestige score (an O*NET-derived or Hauser-Warren-style socioeconomic index, SEI); $d{=}4$ an earnings-growth slope rather than a level (non-wage rewards, $d{=}5$, are the one coordinate even a resume corpus cannot directly supply, so the constructive fix targets $d\in\{1,2,3,4,6\}$); $d{=}6$ the median wage we already use. The estimator is deliberately the \emph{same} rank-gap functional as the paper's headline, applied coordinate by coordinate, because keeping the functional form fixed is what makes the dimensions comparable: any difference across $d$ is then attributable to the ER coordinate, not to a change of method. What this buys is exactly the discrimination the residual-as-divergence reading lacks. If $\mathrm{gap}^{(6)}_f$ (wages) is large but $\mathrm{gap}^{(1)}_f$, $\mathrm{gap}^{(3)}_f$ (employer-selectivity, occupational status) are small, the field's apparent ``divergence'' was a \emph{wage-measurement} phenomenon --- prestige does track where graduates land, just not what they are paid. If instead the gap is large \emph{across all $d$}, the decoupling is a \emph{general} labour-market-reward phenomenon and the divergence reading earns its keep. The per-dimension gap thus separates ``wages do not price this field's reward'' from ``the academy and the market genuinely disagree'' --- the very confound this section identified --- which a single scalar can never do. What it assumes is only that the additional coordinates can be measured at the institution-field grain with coverage comparable to earnings; that assumption is what the Revelio-style resume corpus is meant to discharge.

Second, and complementarily, one can collapse the coordinates back into a single \emph{ER index} (a weighted or factor-analytic composite of employer selectivity, placement, occupational status, trajectory and wage) and recompute the headline gap against that index, testing whether the integrated-decoupled ordering of \S4 is robust to enriching ER beyond the wage. The index route answers ``is the typology stable under a richer ER?''; the per-dimension route answers ``is the typology a wage artifact or a general reward phenomenon?''. Together they would convert the present, honestly-hedged claim --- a measurable wedge between prestige and \emph{pay} --- into the stronger claim the title gestures at --- a wedge between prestige and labour-market \emph{standing}. Until that data is in hand, the asymmetry stands as the binding interpretive limit of the design: we have measured one side of the wedge as a network and the other as a number, and the residual we cannot explain is, in part, the price of that imbalance rather than a discovery about the world.% =====================================================================
\section{What the static paper establishes, and what remains open --- all on the ER side}\label{sec:closing}

It is worth closing the retrospective by drawing a single sharp line between what the static analysis has
\emph{earned} and what it has only \emph{gestured at}, because the two fall on opposite sides of the
measurement asymmetry developed in \S\ref{sec:asymmetry}.

\paragraph{Established, and robustly.} The construction earns the following, in roughly descending order of
confidence. (i) The gap is \emph{real signal, not an artifact}: most observed gaps exceed a calibrated
perfect-agreement-plus-noise null ($15\%$ inside the band, $65\%$ above), the global-artifact channel explains
only $R^2=0.20$ of cross-field variance, and what it leaves behind is structured --- subject always to the
per-field reliability filter that admits well-measured high-gap fields (Biology) and screens noise-dominated
ones (Chemistry). (ii) The gap is \emph{structured by discipline}: a precision-weighted multilevel
decomposition attributes $\mathrm{ICC}=0.30$ of cross-field variance to between-CIP-2 clusters, rising to $0.45$
once within-cluster sampling noise is removed and to $0.57$ on the expanded field universe, with a one-way
$\eta^2=0.50$; the integrated--decoupled typology is not an artifact of cluster size (subsampled rank Spearman
$+0.98$). (iii) \emph{One} external channel moves the gap cleanly --- occupational licensing,
$b_{\text{licensure}}=+0.66$ against a near-orthogonal absorption null --- and the search for further channels
returns collinear redundancy ($R^2$ $0.27\!\to\!0.29$, residual ranking unchanged at Spearman $+1.00$), leaving
a large irreducible residual. (iv) The structure is \emph{robust to a completely independent earnings source}:
PSEO UI wages and Scorecard Title-IV medians agree on institution pay at Spearman $+0.94$ and reproduce the
field-level gap at $+0.68$, with the rank gap surviving the netting of destination geography ($0.68\!\to\!0.67$)
even though geography drives earnings \emph{levels}. (v) Three intuitive alternative readings are \emph{ruled
out} as scope conditions: the gap is not generic-brand-versus-field-specific prestige on medians, not a distinct
``shape'' beyond its strength, and program-level ``mispricing'' is mostly geography and selection. (vi) The
labour market \emph{actively reprices fields} in real dollars over two decades ($46/60$ fields rising;
Statistics $+\$2{,}311$/yr, Computer Science $+\$1{,}616$/yr; Pharmacy $-\$1{,}774$/yr). Each of these is a
statement about \emph{prestige versus the median wage}, and each is measured on the well-instrumented side or
cross-validated on it.

\paragraph{Open, and concentrated on the ER side.} Everything the analysis cannot yet settle lives on the
employer-reputation axis, and \S\ref{sec:asymmetry} explains why this is no coincidence. The two-signal model's
most ambitious claim --- that the licensing-netted residual estimates valuation divergence $(1-\rho_f)$ --- is
\textsc{unconfirmed}: its make-or-break external test failed, with two divergence proxies of \emph{opposite}
sign (O*NET $+0.41$, SDR $-0.61$) that are themselves negatively correlated ($-0.64$), exactly the signature
expected when there is no single well-defined divergence object for a scalar residual to estimate. The lead--lag
direction between prestige and earnings is \emph{not identified} (two noisy periods; prestige stability
$\rho\approx0.50$ is only suggestive of a slow anchor). And the causal, tail-resolved payoff of elite brand that
\citet{chetty2023diversifying} locate above the median is, by construction, invisible to a median-based design.
None of these is a failure of effort; each is a failure of \emph{measurement reach} on the ER axis.

\paragraph{The one closing sentence.} The static paper has measured, validated, and structured a wedge between
academic prestige and the \emph{median wage}, and has been honest that ``employer reputation'' is the broader
construct the wage proxies for; the binding constraint on every remaining question is the measurement asymmetry,
and its resolution is the multi-dimensional ER of \S\ref{sec:asymmetry} --- the data the static design does not
yet have, and the natural next study.

% =====================================================================
\section*{Reproducibility and provenance}
Every statistic in this companion is copied verbatim from a \texttt{*\_RESULT.md} file in the analysis record
and indexed, claim by claim, in \texttt{paper/CLAIMS\_LEDGER.md}; no number here was recomputed or invented for
the retrospective. The estimators and design rationales are reconstructed from those same RESULT files and from
the generating scripts they name. The two documents share a preamble and bibliography (\texttt{refs.bib}) and are
intended to compile in one style: \texttt{pdflatex methods\_companion} $\to$ \texttt{bibtex methods\_companion}
$\to$ \texttt{pdflatex} ($\times 2$).

\bibliographystyle{plainnat}
\bibliography{refs}

\end{document}
